%%%%%%%%%%%%%%%%%%%%%%%%%%%%%%%%%%%%%%%%%
% Beamer Presentation
% LaTeX Template
% Version 1.0 (10/11/12)
%
% This template has been downloaded from:
% http://www.LaTeXTemplates.com
%
% License:
% CC BY-NC-SA 3.0 (http://creativecommons.org/licenses/by-nc-sa/3.0/)
%
%%%%%%%%%%%%%%%%%%%%%%%%%%%%%%%%%%%%%%%%%

%----------------------------------------------------------------------------------------
%	PACKAGES AND THEMES
%----------------------------------------------------------------------------------------

\documentclass[10pt]{beamer}

\mode<presentation>{

% The Beamer class comes with a number of default slide themes
% which change the colors and layouts of slides. Below this is a list
% of all the themes, uncomment each in turn to see what they look like.

%\usetheme{default}
%\usetheme{AnnArbor}
\usetheme{Antibes}  %Good
%\usetheme{Bergen}
%\usetheme{Berkeley}
%\usetheme{Berlin}
%\usetheme{Boadilla}  %Good
%\usetheme{CambridgeUS}
%\usetheme{Copenhagen}
%\usetheme{Darmstadt}
%\usetheme{Dresden}
%\usetheme{Frankfurt}
%\usetheme{Goettingen}
%\usetheme{Hannover}
%\usetheme{Ilmenau}
%\usetheme{JuanLesPins}
%\usetheme{Luebeck}
%\usetheme{Madrid}  %Good
%\usetheme{Malmoe}
%\usetheme{Marburg}
%\usetheme{Montpellier}
%\usetheme{PaloAlto}
%\usetheme{Pittsburgh}
%\usetheme{Rochester}  %Good
%\usetheme{Singapore}  %Good
%\usetheme{Szeged}
%\usetheme{Warsaw}

% As well as themes, the Beamer class has a number of color themes
% for any slide theme. Uncomment each of these in turn to see how it
% changes the colors of your current slide theme.

%\usecolortheme{albatross}
%\usecolortheme{beaver}
%\usecolortheme{beetle}
%\usecolortheme{crane}
\usecolortheme{dolphin}  %Good
%\usecolortheme{dove}
%\usecolortheme{fly}
%\usecolortheme{lily}
%\usecolortheme{orchid}  %Good
%\usecolortheme{rose}  %Good
%\usecolortheme{seagull}
%\usecolortheme{seahorse}  %Good
%\usecolortheme{whale}  %Good
%\usecolortheme{wolverine}

%\setbeamertemplate{footline} % To remove the footer line in all slides uncomment this line
%\setbeamertemplate{footline}[page number] % To replace the footer line in all slides with a simple slide count uncomment this line

%\setbeamertemplate{navigation symbols}{} % To remove the navigation symbols from the bottom of all slides uncomment this line
}

\usepackage{graphicx} % Allows including images
\usepackage{booktabs} % Allows the use of \toprule, \midrule and \bottomrule in tables
\usepackage{amsmath}
\usepackage{listings}
\usepackage{courier}           % for the courier font (optional)
\usepackage{multicol}          % for two equations side by side
\usepackage{centernot}
\usepackage{hyperref}
\usepackage{xcolor}
\usepackage{stmaryrd}
\usepackage{hieroglf}
\usepackage{scalerel}
\usepackage{enumerate}
\usepackage{courier}
\usepackage{appendixnumberbeamer}
\usepackage{pgfplots}
\usepackage{mdframed} %nice frames
\usepackage{bm}
\usepackage{multirow}
\usepackage{tikz-cd}

\hypersetup{ % play with these to change the look of hyperlinks
	colorlinks=true,
%	linkcolor=white,
%	filecolor=white,
  allcolors=.,
	urlcolor=blue,
%	citecolor=black
}


%\setbeamertemplate{bibliography item}{\insertbiblabel}


%------------------------------------------------------------------------------------------
%    ENVIRONMENT SETTINGS
%------------------------------------------------------------------------------------------

% TEMPLATE SETTINGS
\setbeamertemplate{theorems}[numbered]
\setbeamertemplate{sidebar right}{}
\setbeamertemplate{footline}{%
	\hspace{1em}\small\insertframenumber/\inserttotalframenumber
	\hfill\usebeamertemplate***{navigation symbols}}

\makeatletter
\renewcommand{\itemize}[1][]{%
  \beamer@ifempty{#1}{}{\def\beamer@defaultospec{#1}}%
  \ifnum \@itemdepth >2\relax\@toodeep\else
    \advance\@itemdepth\@ne
    \beamer@computepref\@itemdepth% sets \beameritemnestingprefix
    \usebeamerfont{itemize/enumerate \beameritemnestingprefix body}%
    \usebeamercolor[fg]{itemize/enumerate \beameritemnestingprefix body}%
    \usebeamertemplate{itemize/enumerate \beameritemnestingprefix body begin}%
    \list
      {\usebeamertemplate{itemize \beameritemnestingprefix item}}
      {\setlength\itemsep{1em}% NEW
        \def\makelabel##1{%
          {%
            \hss\llap{{%
                \usebeamerfont*{itemize \beameritemnestingprefix item}%
                \usebeamercolor[fg]{itemize \beameritemnestingprefix item}##1}}%
          }%
        }%
      }
  \fi%
  \beamer@cramped%
  \raggedright%
  %\beamer@firstlineitemizeunskip%
}
\makeatother

% ENUMERATE COLOR SETTINGS
\newenvironment{doneenv}{\only{\setbeamercolor{local structure}{fg=gray}\color{gray}}}{}
\newenvironment{comeenv}{\only{\setbeamercolor{local structure}{fg=red}\color{red}}}{}

% LSTLISTING SETTINGS
\lstdefinestyle{myStyle}{
	%	language=Coq,
	keywords={Inductive,Require,Import,Definition,Fixpoint,match,with,end,let,in,fix},
	basicstyle=\footnotesize\ttfamily,
	keywordstyle=\color{green}, % Blue clashes with the cyan links. Change if you want.
	stepnumber=1,
	tabsize=2,
	numbers=none,
	numberstyle=\tiny,
	numbersep=5pt,
	showspaces=false,
	escapechar=`,
	showstringspaces=false,
	backgroundcolor = \color{lightgray},
}

\lstset{basicstyle=\ttfamily\footnotesize,breaklines=true}
%basicstyle=\fontsize{10}{11}\selectfont\ttfamily,

\lstdefinestyle{myTinyStyle}{
	%   language=Coq,
	keywords={Inductive,Require,Import,Definition,Fixpoint,match,with,end,let,in,fix},
	basicstyle=\normalfont\fontsize{8}{8.3}\tt,
	keywordstyle=\color{green}, % Blue clashes with the cyan links. Change if you want.
	stepnumber=1,
	tabsize=2,
	numbers=none,
	numberstyle=\tiny,
	numbersep=5pt,
	showspaces=false,
	escapechar=`,
	showstringspaces=false
}

\lstset{style=myStyle}
\makeatletter
\newlength{\@mli}
\newcommand{\mli}[1]{%
	\settowidth{\@mli}{\lstinline/#1/}
	\hspace{-.5ex}\begin{minipage}[t]{\@mli}\lstinline/#1/\end{minipage}}
\makeatother
\newcommand{\li}[1]{\ifmmode\mbox{\mli{#1}}\else\mbox{\lstinline/#1/}\fi}

\newcommand{\coq}{\scalebox{.6}{\textpmhg{\Ha}}}

\newcommand{\cdt}[2]{#1\mathrel{\mathchoice{\mkern-5mu}{\mkern-5mu}{\mkern-2mu}{\mkern-2mu}} ^{\mathcal{C}}_{#2} \mathrel{\mathchoice{\mkern-3.5mu}{\mkern-3.5mu}{\mkern-2mu}{\mkern-2mu}}}

\newcommand{\cdw}[2]{#1\mathrel{\mathchoice{\mkern-5mu}{\mkern-5mu}{\mkern-2mu}{\mkern-2mu}} ^{\mathcal{CW}}_{#2} \mathrel{\mathchoice{\mkern-3.5mu}{\mkern-3.5mu}{\mkern-2mu}{\mkern-2mu}}}

\newcommand{\cds}[2]{#1\mathrel{\mathchoice{\mkern-5mu}{\mkern-5mu}{\mkern-2mu}{\mkern-2mu}} ^{\mathcal{CS}}_{#2} \mathrel{\mathchoice{\mkern-3.5mu}{\mkern-3.5mu}{\mkern-2mu}{\mkern-2mu}}}

\newcommand{\rf}[2]{{#1\mathrel{\mathchoice{\mkern-5mu}{\mkern-5mu}{\mkern-2mu}{\mkern-2mu}} ^{\mathcal{R}}_{#2} \mathrel{\mathchoice{\mkern-3.5mu}{\mkern-3.5mu}{\mkern-2mu}{\mkern-2mu}}}}

\newcommand{\invhyper}{I}
\renewcommand{\angle}[1]{\langle #1\rangle}

\newcommand{\W}{\mathcal{W}}
\newcommand{\contract}{\text{\scshape{cont}}}
\newcommand{\repeatable}{\text{\scshape{rept}}}

\newcommand{\divc}{\text{divc}}
\newcommand{\logc}{\text{logc}}
\newcommand{\logstarc}{\text{logc*}}

\newcommand{\runtime}{\mathcal{T}}
\newcommand{\Texp}{\runtime_{\li{exp}}}
\newcommand{\Tmul}{\runtime_{\li{mul}}}
\newcommand{\Tleb}{\runtime_{\li{leb}}}
\newcommand{\Tsucc}{\runtime_{\li{succ}}}
\newcommand{\bigO}{\text{O}}

\newcommand{\Ack}{\ensuremath{\mathcal{A}}}
\newcommand{\CA}{\ensuremath{\text{C}}}
%\renewcommand{\S}{\ensuremath{\text{S}}}
\newcommand{\CR}{\ensuremath{\text{CR}}}
\newcommand{\CRH}{\ensuremath{\text{CRH}}}
\newcommand{\wt}[1]{\ensuremath{\widetilde{#1}}}
\newcommand{\IAR}{\ensuremath{\text{IAR}}}
\newcommand{\IARH}{\ensuremath{\text{IARH}}}

\newcommand\hide[1]{}
\newenvironment{centermath}
{\begin{center}$\displaystyle}
	{$\end{center}}
\renewcommand{\note}[2][polish]{{\color{red} #2}{\marginpar{\tiny \color{blue} #1}}}
\renewcommand{\implies}{\Rightarrow}
\renewcommand{\iff}{\Leftrightarrow}


% TABLE, ARRAY SETTINGS
\newcommand{\topspace}[1]{\rule{0pt}{2.6ex+#1}}        % = `top' strut
\newcommand{\botspace}[1]{\rule[-0.9ex-#1]{0pt}{0pt}}   % = `bottom' strut


% DARKENED COLORS
\colorlet{red}{red!80!black}
\colorlet{blue}{blue!80!black}
\colorlet{green}{green!50!black}
\colorlet{lightgray}{gray!30!white}

% FBOX SETTINGS
\newcommand{\imppar}[1]{{\setlength{\fboxrule}{1pt}\fcolorbox{red}{white}{#1}}}
\newcommand{\impinline}[1]{{\setlength{\fboxrule}{1pt}\setlength{\fboxsep}{2pt}\fcolorbox{white}{lightgray}{#1}}}

%\AtBeginSection[]{
%	\begin{frame}
%	\vfill
%	\centering
%	\begin{beamercolorbox}[sep=8pt,center,shadow=true,rounded=true]{title}
%		\Huge{\insertsectionhead}\par%
%	\end{beamercolorbox}
%	\vfill
%	\end{frame}
%}

%\AtBeginSection[*]{
%	\begin{frame}
%	\vfill
%	\centering
%	\begin{beamercolorbox}[sep=8pt,center,shadow=true,rounded=true]{title}
%		\Huge{\insertsectionhead}\par%
%	\end{beamercolorbox}
%	\vfill
%\end{frame}
%}
%----------------------------------------------------------------------------------------
%	TITLE PAGE
%----------------------------------------------------------------------------------------

\title[Inverting the Ackermann Hierarchy]{A Functional Proof Pearl: \protect\\ Inverting the Ackermann Hierarchy} % The short title appears at the bottom of every slide, the full title is only on the title page
\author{\underline{Linh Tran}, Anshuman Mohan, Aquinas Hobor} % Your name
\institute[Yale,NUS]{$
\includegraphics[scale=0.05]{yale-nus-logo.jpg} \quad \quad
\includegraphics[scale=0.4]{NUS_logo_full-horizontal.jpg}\quad \quad
 \includegraphics[scale=0.015]{yale-logo.png}$}
%{$\includegraphics[scale=0.33]{NUS_logo_full-horizontal.jpg}} % Your institution as it will appear on the bottom of every slide, may be shorthand to save space
% {
% National University of Singapore % Your institution for the title page
% }
\date{The 9th ACM SIGPLAN International Conference on Certified Programs and Proofs\\ \today} % Date, can be changed to a custom date

\begin{document}

\begin{frame}[plain]
\titlepage % Print the title page as the first slide
\end{frame}

%----------------------------------------------------------------------------------------
%	PRESENTATION SLIDES
%----------------------------------------------------------------------------------------
% ============================================================================
% OVERVIEW
% ============================================================================

\section*{Introduction}
\label{sec: overview}
\begin{frame}
\frametitle{The Ackermann Function}
	
	The Ackermann-P\'eter function is defined as:
	\begin{equation*}
	A(n, m) \triangleq \begin{cases}
	m + 1 & \text{ when } n = 0 \\
	A(n-1, 1) & \text{ when } n > 0, m = 0 \\
	A\big(n-1, A(n, m-1)\big) & \text{ otherwise}
	\end{cases}
	\end{equation*}
	
	\pause 
	The \emph{diagonal} Ackermann function is $\Ack(n)~\triangleq~A(n, n)$.
	
	\bigskip
	
	\pause 
	First few values of $\Ack(n)$:
	
  \begin{minipage}{0.5\linewidth}
		\begin{equation*}
	\begin{array}{r|ccccc}
	 n & 0 & 1 & 2 & 3 & 4 \\ \hline
	 \Ack(n) & 1 & 3 & 7 & 61 & 2^{2^{2^{65536}}} - 3 \topspace{3pt}
	\end{array}
	\end{equation*}
  \end{minipage}
  \quad \pause 
  \begin{minipage}{0.4\linewidth}
  	\textcolor{red}{Explosive growth!}
  \end{minipage}

\end{frame}


\begin{frame}
\frametitle{The Inverse Ackermann Function}

The \emph{inverse Ackermann function}
%, $\alpha$, maps $n$ to the smallest~$k$ for
%which~$n \le \Ack(k)$, \emph{i.e}
\impinline{$\alpha(n) \triangleq \min\left\{k\in \mathbb{N} : n \le \Ack(k)\right\}$}.

%\smallskip
%
%\pause 
%$\alpha(n)$ grows slowly but is hard to \emph{compute} for large $n$
%\\ because it is entangled with the explosively-growing $\Ack(k)$.

\bigskip

\pause 
\textbf{Naive Approach:} Compute $\Ack(0), \Ack(1), \ldots$ until $n\le \Ack(k)$. Return $k$.

\bigskip

\textbf{Time complexity:} $\Omega(\Ack(\alpha(n)))$.
\\ Computing $\alpha(100) \mapsto^{*} 4$ requires at least
%$\Ack(4) = 2^{2^{2^{65536}}}$ steps!
$\Ack(4) = 2^{2^{2^{65536}}} - 3$ steps!

\bigskip

\pause 
\textbf{Engineering hack:} Hardcode with lookup tables.
\\ $n > 61 \implies \text{ans} = 4$. \textcolor{red}{Wrong for large enough inputs.}

\bigskip

\pause 
\imppar{\text{\textbf{Our Goal.} Compute $\alpha$ for \emph{all inputs} without computing $\Ack$.}}
\end{frame}

\begin{frame}[fragile]
\frametitle{Our Solution}

\vspace{-1em}
\lstset{style=myTinyStyle}
% Linked by A
\begin{mdframed}[backgroundcolor=lightgray, roundcorner=10pt,leftmargin=0, rightmargin=0, innerleftmargin=0, innertopmargin=-5,innerbottommargin=-5, outerlinewidth=0, linecolor=lightgray]
\begin{lstlisting}
Require Import Omega Program.Basics.

`\href{https://github.com/inv-ack/inv-ack/blob/7270e64a2600b771f2b1b1b151f7d13fb2ae6c97/inv_ack_standalone.v#L6-L11}{Fixpoint cdn\_wkr}` (a : nat) (f : nat -> nat) (n b : nat) :=
 match b with 0 => 0 | S b' =>
  if (n <=? a) then 0 else S (cdn_wkr f a (f n) k')
 end.

`\href{https://github.com/inv-ack/inv-ack/blob/7270e64a2600b771f2b1b1b151f7d13fb2ae6c97/inv_ack_standalone.v#L14}{Definition countdown\_to}` a f n := cdn_wkr a f n n.

`\href{https://github.com/inv-ack/inv-ack/blob/7270e64a2600b771f2b1b1b151f7d13fb2ae6c97/inv_ack_standalone.v#L32-L38} {Fixpoint inv\_ack\_wkr}` (f : nat -> nat) (n k b : nat) :=
 match b with 0 => 0 | S b' =>
  if (n <=? k) then k else let g := (countdown_to f 1) in
                      inv_ack_wkr (compose g f) (g n) (S k) b
 end.

`\href{https://github.com/inv-ack/inv-ack/blob/7270e64a2600b771f2b1b1b151f7d13fb2ae6c97/inv_ack_standalone.v#L42-L46}{Definition inv\_ack\_linear}` (n : nat) : nat :=
 match n with 0 | 1 => 0 | _ => 
  let f := (fun x => x - 2) in inv_ack_wkr f (f n) 1 (n - 1)
 end.
\end{lstlisting}
\end{mdframed} 
\end{frame}


%begin{frame}
%\frametitle{Ackermann \emph{vs} Hyperoperation}
%
%The Ackermann function is easy to define, but hard to
%understand.
%
%\bigskip
%
%\pause 	
%Let's index by the first argument. \\ \smallskip
%Define $\Ack_n \triangleq \lambda b.A(n,b)$. \\ \smallskip
%Then, for $n>0$, $\Ack_n$ is the result of applying the previous $\Ack_{n-1}$ $b$ times.
%
%%with a
%%\href{https://github.com/inv-ack/inv-ack/blob/7270e64a2600b771f2b1b1b151f7d13fb2ae6c97/repeater.v\#L161-L177}{\emph{kludge}}. %Linked by Linh
%
%\bigskip
%
%\pause 
%The hierarchical structure resembles that of \textcolor{red}{\emph{hyperoperations}}.
%
%\smallskip
%Studying hyperoperations helps us understand the Ackermann hierarchy.
%
%%To better understand the Ackermann function as a hierarchy and this kludge, we explore the closely-related hyperoperations.
%\end{frame}


\begin{frame}
\frametitle{Introduction: Ackermann \emph{vs} Hyperoperation}

Treating $b$ as the main argument reveals similarities between $A(n, b)$ and the \textcolor{red}{\emph{hyperoperations}} $a[n]b$, while allows the notion of \textcolor{red}{\emph{inverse functions}}.

\begin{table}[t]
	\begin{centermath}
		\begin{array}{c@{\hskip 0.5em}|@{\hskip 1em}c@{\hskip 1em}c@{\hskip 1em}|@{\hskip 1em}c@{\hskip 1em}c}
%			  & \multicolumn{2}{|@{\hskip 0.5em}c@{\hskip 0.5em}|}{\text{Main hierarchies}} & \multicolumn{2}{|@{\hskip 0.5em}c@{\hskip 0.5em}|}{\text{Inverses hierarchies}} \\
			n & a [n] b & A(n, b) & a \angle{n} b & \alpha_n(b)\\
			\hline
			0 & 1 + b & 1 + b & b - 1 & b - 1 \\
			1 & a + b & 2 + b & b - a & b - 2 \\
			2 & a \cdot b & 2b + 3 & \left\lceil \frac{b}{a} \right\rceil & \left\lceil \frac{b-3}{2} \right\rceil \\
			3 & a^b & 2^{b + 3} - 3 & \left\lceil \log_a ~ b \right\rceil & \left\lceil \log_2 ~ (b + 3)\right\rceil - 3 \\
			[1pt]
			4 & \underbrace{a^{.^{.^{.^a}}}}_b & \underbrace{2^{.^{.^{.^2}}}}_{b+3} - 3 & \log^*_a ~ b & \log^*_2 ~ (b + 3) - 3
		\end{array}
	\end{centermath}
	\label{table: hyperop-ack-inv}
\end{table}

\pause
\href{https://github.com/inv-ack/inv-ack/blob/7270e64a2600b771f2b1b1b151f7d13fb2ae6c97/repeater.v\#L161-L177}{Connection:}
%\pause
$A(n, b) = {\color{red}2}[n](b{\color{red}+3}) {\color{red}- 3}$
\quad and \quad $\alpha_n(b) = {\color{red}2}\angle{n}(b{\color{red}+3}) {\color{red}- 3}$.

\end{frame}


\begin{frame}
\frametitle{Introduction: Inverse Hierarchies to Inverse Ackermann}

We explore the upper inverse relation:
\begin{equation*}
\begin{cases}
\forall b. \forall c.\quad b \le \Ack_n(c) & \!\! \iff \ \ \alpha_n(b)\le c \\
\forall b. \forall c. \quad b \le a[n]c & \!\! \iff \ \ a\angle{n}b \le c
\end{cases}
\end{equation*}

\textbf{Redefine $\bm{\alpha}$:}
$\alpha(n) = \min\{k: n\le \Ack_k(k) \} = \min\{k: \alpha_k(n)\le k \}$.

\bigskip

\pause
\textbf{Computing $\bm{\alpha} $ through $\bm{\alpha_i}$!} No need to go through $\Ack$.

\bigskip

\pause
\imppar{\textbf{Goal.} Build the inverse towers independent from the original towers.}

\end{frame}






%\begin{frame}
%\frametitle{}
%\end{frame}

%TODO: think of a better way to make the roadmap i.e. diagram or something

\section*{Roadmap}
\begin{frame}
\frametitle{Roadmap}
\setlength{\leftmargini}{3.8em}
\setbeamertemplate{enumerate items}[default]
\setbeamercolor{local structure}{fg=black}

\imppar{\textbf{Goal. Inverting $\Ack$ - without computing $\Ack$.}}

\bigskip

\begin{enumerate}[\bfseries Step 1.]
  \itemsep 2ex
  \pause
	\item<come@6>
	Build the hyperoperations/Ackermann hierarchies via \textbf{Repeater}.
	
	\item<3->
	Build inverses of the hyperoperations/Ackermann hierarchies \\ via \textbf{Countdown}, the inverse of \emph{Repeater}.
	
	\item<4->
	Use the hierarchy to implement the inverse Ackermann function \\ for inputs encoded in unary. $\bigO (n)$.
\end{enumerate}
\end{frame}


\section{Step 1: Hyperoperations and Ackermann via Repeater}
\label{sec: countdown-repeater}
\subsection{Repeated Application}


\begin{frame}
\frametitle{Hierarchical Relation from Recursive Rule}

Let's look at the recursive rules for the Ackermann function and the hyperoperations.
\begin{equation*}
\begin{array}{lrcl}
\textit{Hyperoperations: } \quad & \textcolor{red}{a[n+1]}(b+1) & ~ \triangleq ~ & \textcolor{green}{a[n]}\big(\textcolor{red}{a[n+1]}b\big)
\\[15pt]
\textit{Ackermann function: } \quad & \textcolor{red}{A(n+1,} b+1) & ~ \triangleq ~ & \textcolor{green}{A\big(n,} \textcolor{red}{A(n+1,} b)\big)
\end{array}
\end{equation*}

\pause
Indexing \impinline{$\Ack_n \triangleq \lambda b. A(n, b)$}: \imppar{$\textcolor{red}{\Ack_{n+1}}(b+1) ~ = ~  \textcolor{green}{\Ack_n}\big(\textcolor{red}{\Ack_{n+1}}(b)\big)$}.

\bigskip

\pause
\emph{The next level can be built from the previous level!}

\end{frame}


\begin{frame}
\frametitle{Building the Next Level}

	\begin{centermath}
		\begin{array}{l@{\hskip 1em}l@{\hskip 1em}|@{\hskip 1em}l}
%			  & \multicolumn{2}{|@{\hskip 0.5em}c@{\hskip 0.5em}|}{\text{Main hierarchies}} & \multicolumn{2}{|@{\hskip 0.5em}c@{\hskip 0.5em}|}{\text{Inverses hierarchies}} \\
		&	\textcolor{red}{a[n+1]}b & \textcolor{red}{\Ack_{n+1}}(b)
		\\[5pt]
	\pause = & \textcolor{green}{a[n]}\big(\textcolor{red}{a[n+1]}(b-1)\big) & \textcolor{green}{\Ack_n}\big(\textcolor{red}{\Ack_{n+1}}(b-1)\big)
	\\[5pt]
	\pause = & \big(\textcolor{green}{a[n]}\circ \textcolor{green}{a[n]}\big)\big(\textcolor{red}{a[n+1]}(b-2)\big) & \big(\textcolor{green}{\Ack_n} \circ \textcolor{green}{\Ack_n}\big)\big(\textcolor{red}{\Ack_{n+1}}(b-2)\big)
	\\[5pt]
	\pause = & \ldots & \ldots
	\\[5pt]
	\pause = & \underbrace{\big( \textcolor{green}{a[n]} \circ \cdots \circ \textcolor{green}{a[n]} \big)}_{b \text{ times}} \big(\textcolor{red}{a[n+1]}0\big) &
	\underbrace{\big( \textcolor{green}{\Ack_n} \circ \cdots \circ \textcolor{green}{\Ack_n} \big)}_{b \text{ times}} \big(\textcolor{red}{\Ack_{n+1}}(0)\big)
	\\[25pt]
	\pause = & \underbrace{\big(\textcolor{green}{a[n]}\big)^{(b)}}_{\substack{\text{repeated} \\ \text{applications}}}
	\underbrace{\big(\textcolor{red}{a[n+1]}0\big)}_{\text{initial value}}
	& \underbrace{\big(\textcolor{green}{\Ack_n}\big)^{(b)}}_{\substack{\text{repeated} \\ \text{applications}}}
	\underbrace{\big(\textcolor{red}{\Ack_{n+1}}(0)\big)}_{\text{initial value}}
		\end{array}
	\end{centermath}

\end{frame}


%\begin{frame}
%\frametitle{Formal Definition of Hyperoperation Hierarchy}
%
%\vspace{-1.5ex}
%\begin{equation*}
%\begin{array}{lrcl}
%\textit{1. 0$^{\textit{th}}$ level: } & \uncover<2->{a[0]b & ~ \triangleq ~ & b + 1} \\[3pt]
%\textit{2. Initial values: } & \uncover<3->{a[n+1]0 & ~ \triangleq ~ &
%\small \begin{cases}
%a & \text{when } n = 0 \\
%0 & \text{when } n = 1 \\
%1 & \text{otherwise}
%\end{cases}} \\[15pt]
%\textit{3. Recursive rule: } \quad & \uncover<4->{a[n+1](b+1) & ~ \triangleq ~ & a[n]\big(a[n+1]b\big)}
%\end{array}
%\end{equation*}
%
%\vspace{-1ex}
%
%\begin{equation*}
%\begin{aligned}
%& a[n+1]b \\
%& \uncover<5->{= ~ a[n]\big(a[n+1](b-1)\big) ~ = ~ a[n]\big(a[n](a[n+1](b-2))\big)} \\
%& \uncover<6->{= ~ \big(a[n]\circ a[n]\big)\big(a[n+1](b-2)\big) ~ = ~ \ldots} \\
%& \uncover<7->{= ~ \underbrace{\big( a[n]\circ a[n]\circ \cdots \circ a[n] \big)}_{b \text{ times}} \big(a[n+1]0\big)  ~ = ~ 
%\underbrace{\big(a[n]\big)^{(b)}}_{\substack{\text{repeated} \\ \text{applications}}}
%\underbrace{\big(a[n+1]0\big)}_{\text{initial value}}}
%\end{aligned}
%\end{equation*}
%
%\end{frame}
%
%
%\begin{frame}
%\frametitle{Formal Definition of Ackermann Hierarchy}
%\begin{equation*}
%\begin{array}{lrcl}
%\textit{1. 0$^{\textit{th}}$ level: } & \uncover<2->{\Ack_0(b) & ~ \triangleq ~ & b + 1} \\[5pt]
%\textit{2. Initial values: } & \uncover<2->{\Ack_{n+1}(0) & ~ \triangleq ~ & \Ack_n(1)} \\[5pt]
%\textit{3. Recursive rule: } \quad & \uncover<2->{\Ack_{n+1}(b+1) & ~ \triangleq ~ & \Ack_n\big(\Ack_{n+1}(b)\big)}
%\end{array}
%\end{equation*}
%
%\bigskip
%
%\begin{equation*}
%\begin{aligned}
%& \Ack_{n+1}(b) \\
%& \uncover<3->{= ~ \Ack_n\big(\Ack_{n+1}(b-1)\big) ~ = ~ \Ack_n\big(\Ack_n\big(\Ack_{n+1}(b-2)\big)\big) \\
%& = ~ \big(\Ack_n\circ \Ack_n\big) \big(\Ack_{n+1}(b-2)\big) ~ = ~ \ldots \\
%& = ~ \underbrace{\big( \Ack_n\circ \Ack_n\circ \cdots \circ \Ack_n \big)}_{b \text{ times}} \big(\Ack_{n+1}(0)\big)  ~ = ~ 
%\underbrace{\big(\Ack_n\big)^{(b)}}_{\substack{\text{repeated} \\ \text{applications}}}
%\underbrace{\big(\Ack_{n+1}(0)\big)}_{\text{initial value}}}
%\end{aligned}
%\end{equation*}
%
%\end{frame}


\subsection{Introducing Repeater}


\begin{frame}[fragile]
\frametitle{Mechanizing Our Observations}
%The next level in the hyperoperations/Ackermann hierarchy is the result of $b$ compositional applications of the current level to an initial value.
%
%\bigskip
%
%We can abstract the concept of repeated application in a higher-order function called \emph{repeater}:
%
%\smallskip

%Writing $a[n+1]b = \big(a[n]\big)^{(b)}\big(a[n+1]0\big)$ or $\Ack_{n+1}(b) = \big(\Ack_n\big)^{(b)}\big(\Ack_{n+1}(0)\big)$ is clumsy since $b$ is no longer the main argument.

% \vspace{-2ex}
% \begin{equation*}
% \begin{array}{l|l|l}
% \text{Expression} & \text{Repeated Applications} & \text{CSL } \texttt{iter} \\ \hline
% a[n+1]b & \big(a[n]\big)^{(b)}\big(a[n+1]0\big) & \texttt{iter}\big(b, a[n], a[n+1]0\big) \topspace{2pt} \botspace{2pt} \\
% \Ack_{n+1}(b) & \big(\Ack_n\big)^{(b)}\big(\Ack_{n+1}(0)\big) & \texttt{iter}\big(b, \Ack_n, \Ack_{n+1}(0)\big) \topspace{2pt}
% \end{array}
% \end{equation*}
% \\ \vspace{-0.5ex}
% The writing is clumsy since $b$ is no longer the main argument.

% \smallskip

% \pause 
%\textbf{Repeater.} A rearranged form of \texttt{iter}, where $b$ becomes the main argument:
%\impinline{$\rf{f}{u}(b) \triangleq \texttt{iter}(b, f, u) = f^{(b)}(u)$}.
\textbf{Repeater.} %A rearranged \texttt{iter}:
\impinline{$\rf{f}{u}(b)~~=~~f^{(b)}(u)~~\approx~~\texttt{iter}(b, f, u)$}.
\pause \\Read $\rf{f}{a}$ \ as ``the \emph{repeater from} $a$ of $f$''.

%\smallskip

%$\forall a\in \mathbb{N}, f: \mathbb{N}\to \mathbb{N}$, the
%\href{https://github.com/inv-ack/inv-ack/blob/7270e64a2600b771f2b1b1b151f7d13fb2ae6c97/repeater.v\#L32-L36}{\emph{repeater from}}
%$a$ of $f$, denoted by $\rf{f}{a}$ , is a function $\mathbb{N}\to \mathbb{N}$ such that $\rf{f}{a}(n) = f^{(n)}(a)$.

%\pause
\bigskip
% Linked by A
\begin{lstlisting}
`\href{https://github.com/inv-ack/inv-ack/blob/7270e64a2600b771f2b1b1b151f7d13fb2ae6c97/repeater.v#L32-L36}{Fixpoint repeater\_from}` (f : nat -> nat) (a n : nat) : nat :=
  match n with
  | 0 => a
  | S n' => f (repeater_from f a n') end.
\end{lstlisting}

\smallskip

%\begin{minipage}{0.25\linewidth}
%	Drop $b$ to form \\ a recursive rule:
%\end{minipage}
\pause 
\begin{minipage}{0.4\linewidth}
	Drop $b$ to form higher-order \\ recursive rule between levels:
\end{minipage}
\quad 
$\begin{cases}
\displaystyle a[n+1] & = \rf{\big(a[n]\big)}{a[n+1]0} \\[5pt]
\Ack_{n+1} & = \rf{\big(\Ack_n\big)}{\Ack_{n+1}(0)}
\end{cases}$

\end{frame}


\begin{frame}[fragile]
\frametitle{Hyperoperations via Repeater}
Without Repeater (via double recursion):
\begin{lstlisting}
`\href{https://github.com/inv-ack/inv-ack/blob/7270e64a2600b771f2b1b1b151f7d13fb2ae6c97/repeater.v#L51-L52}{Definition hyperop\_init}` (a n : nat) : nat :=
  match n with 0 => a | 1 => 0 | _ => 1 end.

`\href{https://github.com/inv-ack/inv-ack/blob/7270e64a2600b771f2b1b1b151f7d13fb2ae6c97/repeater.v#L55-L64}{Fixpoint hyperop\_original}` (a n b : nat) : nat :=
  match n with
  | 0    => 1 + b
  | S n' => let fix hyperop' (b : nat) := match b with
              | 0    => hyperop_init a n'
              | S b' => hyperop_original a n' (hyperop' b')
              end in hyperop' b
  end.
\end{lstlisting}

\pause 
With Repeater:
% Linked by Anshuman
\begin{lstlisting} 
`\href{https://github.com/inv-ack/inv-ack/blob/7270e64a2600b771f2b1b1b151f7d13fb2ae6c97/repeater.v#L67-L71}{Fixpoint hyperop}` (a n b : nat) : nat :=
  match n with
  | 0    => 1 + b
  | S n' => repeater_from (hyperop a n') (hyperop_init a n') b
  end.
\end{lstlisting}
\end{frame}


\begin{frame}[fragile]
\frametitle{Ackermann via Repeater}
Without Repeater (via double recursion):
\begin{lstlisting}
`\href{https://github.com/inv-ack/inv-ack/blob/7270e64a2600b771f2b1b1b151f7d13fb2ae6c97/repeater.v#L123-L132}{Fixpoint ackermann\_original}` (m n : nat) : nat :=
  match m with
  | 0    => 1 + n
  | S m' => let fix ackermann' (n : nat) : nat := match n with
              | 0    => ackermann_original m' 1
              | S n' => ackermann_original m' (ackermann' n')
              end in ackermann' n
  end.
\end{lstlisting}

\pause 
With Repeater:
% Linked by A
\begin{lstlisting}
`\href{https://github.com/inv-ack/inv-ack/blob/7270e64a2600b771f2b1b1b151f7d13fb2ae6c97/repeater.v#L135-L139}{Fixpoint ackermann}` (n m : nat) : nat :=
  match n with
  | 0    => S m
  | S n' => repeater_from (ackermann n') (ackermann n' 1) m
  end.
\end{lstlisting}
\end{frame}

%
%\begin{frame}
%\frametitle{}
%\end{frame}


\section*{Roadmap}
\begin{frame}
\frametitle{Roadmap}
\setlength{\leftmargini}{3.8em}
\setbeamertemplate{enumerate items}[default]
\setbeamercolor{local structure}{fg=black}

\imppar{\textbf{Goal.} Inverting $\Ack$ - without computing $\Ack$.}

\bigskip

\begin{enumerate}[\bfseries Step 1.]
	\itemsep 2ex
	\item<done@1->
	Build the hyperoperations/Ackermann hierarchies via \textbf{Repeater}.
	
	\item<come@2->
	Build inverses of the hyperoperations/Ackermann hierarchies via \textbf{Countdown}, the inverse of \emph{Repeater}.
	
	\item
	Use the hierarchy to implement the inverse Ackermann function \\ for inputs encoded in unary. $\bigO (n)$.
\end{enumerate}
\end{frame}


\section{Step 2: Inverting the Hierarchies via Countdown}
\label{sec: countdown}
\subsection{Inverting Repeater with Countdown}

%\begin{frame}
%\frametitle{Repeatable Functions}
%
%Functions in the Ackermann and hyperoperation (when $a\ge 2$) hierarchies are all \textcolor{red}{\emph{repeatable function}}.
%
%\bigskip
%%
%%\pause
%%\textbf{Expansions.} A function $F:\mathbb{N}\to\mathbb{N}$ is an
%%\href{https://github.com/inv-ack/inv-ack/blob/7270e64a2600b771f2b1b1b151f7d13fb2ae6c97/increasing_expanding.v\#L80-L82}{\emph{expansion}} if $\forall n.~ F(n)\ge n$. Further, for $a\in \mathbb{N}$, an expansion $F$ is
%%\href{https://github.com/inv-ack/inv-ack/blob/7270e64a2600b771f2b1b1b151f7d13fb2ae6c97/increasing_expanding.v\#L84-L86}{\emph{strict from}} $a$ if ~$\forall n \ge a.~ F(n) > n$.
%%
%%\bigskip
%
%\pause
%\textbf{Repeatable functions:} A $F$ is \emph{repeatable} from $a$ if\\
%$F$ is strictly increasing and $F$ is an
%\href{https://github.com/inv-ack/inv-ack/blob/7270e64a2600b771f2b1b1b151f7d13fb2ae6c97/increasing_expanding.v\#L80-L82}{\emph{expansion}}
%that is
%\href{https://github.com/inv-ack/inv-ack/blob/7270e64a2600b771f2b1b1b151f7d13fb2ae6c97/increasing_expanding.v\#L84-L86}{\emph{strict from}}
%$a$, i.e. ~$\forall n \ge a.~ F(n) > n$.
%
%\bigskip
%\pause
%We extend our scope of study from functions in the hyperoperations and Ackermann hierarchies to repeatable functions.
%
%\bigskip
%\pause
%\textbf{Advantage.} If $F$ is repeatable from $a$, $F^{-1}$ makes sense and is total,
%\\ and $\rf{F}{a}$ is repeatable from $1$.
%
%\end{frame}
%
%
%
%\begin{frame}
%\frametitle{Inverting Repeater: The Idea of Countdown}
%
%The \href{https://github.com/inv-ack/inv-ack/blob/7270e64a2600b771f2b1b1b151f7d13fb2ae6c97/inverse.v\#L28-L45}{
%	\emph{upper inverse}} of $F$, written $F^{-1}$,
%% {\color{red} $F^{-1}$, $F^{-1}_{\mathit{\shortuparrow}}$, $F^{-1}_{\upharpoonleft}$}
%is \impinline{$\lambda n. \min\{m : F(m)\ge n\}$}.
%
%\smallskip
%
%\pause 
%\textbf{\href{https://github.com/inv-ack/inv-ack/blob/7270e64a2600b771f2b1b1b151f7d13fb2ae6c97/inverse.v\#L65-L77}{Logical equivalence} (more useful)}: If $F:\mathbb{N}\to \mathbb{N}$ is increasing, \\
%then $f = F^{-1}$ iff \impinline{$\forall n, m.~ f(n)\le m \iff n \le F(m)$}.
%
%\bigskip
%
%\textbf{Idea.} Build $\left(\rf{F}{a}\ \right)^{-1}$ from $f \triangleq F^{-1}$.
%	\begin{equation*}
%	\displayindent0pt
%	\displaywidth\textwidth
%	\begin{array}{ccll}
%	& \left(\rf{F}{a}\ \right)^{-1}(n)\le m \pause & \iff \ n \le \rf{F}{a}(m) \pause & \iff n \ \le F^{(m)}(a) \\
%  \pause \iff & f(n)\le F^{(m-1)}(a) \pause & \iff \ \cdots \pause &
%	\iff \ f^{(m)}(n) \le a 
%	\end{array}
%	\end{equation*}
%
%\pause
%\imppar{$\left(\rf{F}{a}\ \right)^{-1}(n)$ is \emph{the least} $m$ for which $f^{(m)}(n)\le a$.}
%
%\end{frame}



\begin{frame}[fragile]

\textbf{Goal.} Construct $\alpha_{n+1}$ from $\alpha_n$ and $a\angle{n+1}$ from $a\angle{n}$.
$$
\tikzset{
LA/.style = {
draw=red, % just to demonstrate, where LA is used
line width=1.5pt, -{Straight Barb[length=3pt]}},
LA/.default=1pt,
LB/.style = {
draw=blue, % just to demonstrate, where LA is used
line width=1.5pt, -{Straight Barb[length=3pt]}},
LB/.default=1pt,
invisible/.style={opacity=0},
visible on/.style={alt={#1{}{invisible}}},
alt/.code args={<#1>#2#3}{%
  \alt<#1>{\pgfkeysalso{#2}}{\pgfkeysalso{#3}}%
  }
}
\newcommand{\ard}{
  \ar[visible on=<2->, d, LB,
    "\ \text{\textbf{\textcolor{blue}{Repeater}}}"]}
\newcommand{\rd}{\ar[d, LA]}
\newcommand{\abd}{
  \ar[visible on=<3->, d, LA,
    "\ \text{\textbf{\textcolor{red}{???}}}"]}
\newcommand{\bd}{\ar[d, LB]}
\begin{tikzcd}[row sep = 13.5pt]
a[n]b & \Ack_n(b) & a\angle{n}b & \alpha_n(b) \\[-15pt] \hline
1 + b \ard & 1 + b \ard & b - 1 \abd & b - 1 \abd
\\
a + b \ard & 2 + b \ard & b - a \abd & b - 2 \abd
\\
a \cdot b \ard & 2b + 3 \ard & \left\lceil \frac{b}{a} \right\rceil \abd & \left\lceil \frac{b-3}{2} \right\rceil \abd
\\
a^b & 2^{b + 3} - 3 & \left\lceil \log_a b \right\rceil & \left\lceil \log_2 (b + 3)\right\rceil - 3
%\\
%\underbrace{a^{.^{.^{.^a}}}}_b & \underbrace{2^{.^{.^{.^2}}}}_{b+3} - 3 & \log^*_a ~ b & \log^*_2 ~ (b + 3) - 3
\end{tikzcd}
$$

\uncover<4->{\imppar{\textbf{Next Goal.} Construct the inverse of $\rf{F}{a}\ $ entirely from $F$'s inverse.}}

\end{frame}


%\begin{frame}
%\frametitle{Formalizing Inverse}
%
%If $F: \mathbb{N} \to \mathbb{N}$ is strictly increasing (holds for all $\Ack_n$ and $a[n]$ when $a\ge 2$),
%the \href{https://github.com/inv-ack/inv-ack/blob/7270e64a2600b771f2b1b1b151f7d13fb2ae6c97/inverse.v\#L28-L45}{
%	\emph{upper inverse}} of $F$, written $F^{-1}$,
%% {\color{red} $F^{-1}$, $F^{-1}_{\mathit{\shortuparrow}}$, $F^{-1}_{\upharpoonleft}$}
%is \impinline{$\lambda n. \min\{m : F(m)\ge n\}$}.
%
%\smallskip
%
%\textbf{Note.} Since $F$ is strictly increasing, $F^{-1}$ is a total function.
%
%\bigskip
%
%\pause 
%\textbf{\href{https://github.com/inv-ack/inv-ack/blob/7270e64a2600b771f2b1b1b151f7d13fb2ae6c97/inverse.v\#L65-L77}{Logical equivalence} (more useful)}:
%\impinline{$\forall n, m.~ F^{-1}(n)\le m \iff n \le F(m)$}.
%
%\bigskip
%
%\pause
%\textbf{Repeatable functions.} $F$ is \textcolor{red}{\emph{repeatable from}} $a$ if $F$ is strictly increasing and \textcolor{red}{\emph{expanding from}} $a$, i.e. $~\forall n.~F(n) \ge n~$ and $~\forall n \ge a.~F(n) > n$.
%
%\vspace{5pt}
%
%\pause
%\textbf{Advantage.} $\displaystyle \begin{cases}
%F \text{ is repeatable from } a \implies \rf{F}{a}\ \text{ is repeatable from } 0. \\[3pt]
%\forall n, \ \Ack_n \text{ is repeatable from } 0. \\[3pt]
%\forall n, \forall a \ge 2, \ a[n] \text{ is repeatable from } a[n+1](0).
%\end{cases}$
%
%\end{frame}


\begin{frame}
\frametitle{Formalizing Inverse}

\textbf{Assumption.} We assume $a\ge 2$ whenever $a[n]b$ is mentioned.

\bigskip

\pause

The \href{https://github.com/inv-ack/inv-ack/blob/7270e64a2600b771f2b1b1b151f7d13fb2ae6c97/inverse.v\#L28-L45}{
	\emph{upper inverse}} of $F$, written $F^{-1}$,
% {\color{red} $F^{-1}$, $F^{-1}_{\mathit{\shortuparrow}}$, $F^{-1}_{\upharpoonleft}$}
is \impinline{$\lambda n. \min\{m : F(m)\ge n\}$}.

\smallskip

\textbf{Remarks.}

$\bullet~$ If $F$ is strictly increasing, $F^{-1}$ is a total function.

$\bullet~$ $\Ack_n$ and $a[n]$ are strictly increasing for each $n$.




\bigskip

\pause 
\textbf{\href{https://github.com/inv-ack/inv-ack/blob/7270e64a2600b771f2b1b1b151f7d13fb2ae6c97/inverse.v\#L65-L77}{Logical equivalence} (more useful)}:
\impinline{$\forall n, m.~ F^{-1}(n)\le m \iff n \le F(m)$}.

\end{frame}



\begin{frame}
\frametitle{Inverting Repeater: The Idea of Countdown}

Assume F satisfies the 3 conditions of the theorem. \\ \smallskip
Then \impinline{$\forall n, m.~ F^{-1}(n)\le m \iff n \le F(m)$}.

\smallskip

Let us build $\left(\rf{F}{a}\ \right)^{-1}$ from $F^{-1}$.
\pause
	\begin{equation*}
	\displayindent0pt
	\displaywidth\textwidth
	\begin{array}{ccll}
	& \left(\rf{F}{a}\ \right)^{-1}(n)\le m \pause & \iff \ n \le \rf{F}{a}\ (m) \pause & \iff n \ \le F^{(m)}(a) \\[5pt]
  \pause \iff & F^{-1}(n)\le F^{(m-1)}(a) \pause & \iff \ \cdots \pause &
	\iff \ \left(F^{-1}\right)^{(m)}(n) \le a 
	\end{array}
	\end{equation*}

\pause
\imppar{$\left(\rf{F}{a}\ \right)^{-1}(n)$ is \emph{the least} $m$ for which $\left(F^{-1}\right)^{(m)}(n)\le a$.}

\end{frame}


%\begin{frame}
%\frametitle{Formalizing countdown}
%
%\textbf{Countdown.} The \textit{\textcolor{red}{countdown} to} $a$ of $f$, written~$\cdt{f}{a}(n)$, if exists, is the smallest number of times $f$ needs to be compositionally applied to
%$n$ for the answer to equal or go below $a$. \emph{i.e.},
%\begin{equation*}
%\cdt{f}{a}(n) \triangleq \min\{m: f^{(m)}(n)\le a \}.
%\end{equation*}
%
%\textbf{Contractions.} $f$ is a \emph{\textcolor{red}{contraction} strict above} a if \\
%$\forall n.~f(n)\le n~$ and $~\forall n > a.~f(n) < n$.
%
%\smallskip
%
%\pause
%If $f$ is a contraction strict from $a$, $\cdt{f}{a}$ is a total function.
%
%\bigskip
%
%\pause
%\textbf{Theorem.} If $F$ is repeatable from $a$,\\
%$F^{-1}$ is a contraction strict above $a$ and $\cdt{\left(F^{-1}\right)}{a}\ \ = \ \left(\rf{F}{a}\ \right)^{-1}$.
%
%\end{frame}


\begin{frame}
\frametitle{Formalizing countdown}

\textbf{Countdown.} The \textit{\textcolor{red}{countdown} to} $a$ of $f$, written~$\cdt{f}{a}(n)$, if exists, is the smallest number of times $f$ needs to be compositionally applied to
$n$ for the answer to equal or go below $a$. \emph{i.e.},
\begin{equation*}
\cdt{f}{a}(n) \triangleq \min\{m: f^{(m)}(n)\le a \}.
\end{equation*}

\pause
\textbf{Theorem.}
$~\displaystyle
\begin{cases}
F \text{ strictly increasing} \\ \forall n.~F(n)\ge n \\
\forall n > a.~F(n) > n
\end{cases}
\quad \pause \implies \quad
\begin{cases}
\cdt{\left(F^{-1}\right)}{a} \text{ is total} \\[3pt]
\cdt{\left(F^{-1}\right)}{a}\ \ = \ \left(\rf{F}{a}\ \right)^{-1}
\end{cases}
$

\bigskip
\pause
\textbf{Corollary.} $a[n+1]$ and $\alpha_{n+1}$ are respectively the\\ countdowns of $a[n]$ and $\alpha_n$ to some appropriate value.

\end{frame}


\begin{frame}[fragile]
\frametitle{A Coq Computation of Countdown}
\textbf{Idea.} %To compute $\cdt{f}{a}(n)$,
Compute $f^{k}(n)$ for $k=0, 1,\ldots$ until $f^{(k)}(n)\le a$. Return $k$.

\bigskip

\pause 
\textbf{Primary Issue.} Termination in Coq.
%Coq needs a known terminating point, i.e. an explicit decreasing argument. How to know when to terminate beforehand?

\bigskip

\textbf{Secondary issue.} $\cdt{f}{a}$ may not exist for general $f$.
%It's hard to define a Coq definition of Countdown for contractions only. Defining a computation for \emph{all functions} that actually computes Countdown on contractions is easier.

\bigskip

\pause 
\textbf{The worker.} $\begin{cases}
\text{Budget } b: & \hspace{-1em}\text{Maximum } b \text{ steps. } \\
\text{Step } i: & \hspace{-1em}\text{Compute } f^{(i)}(n). \\
\text{Stops when: } & \hspace{-1em}\text{budget reaches } 0 \text{ or } f^{(i)}(n) \le a.
\end{cases}$
\end{frame}


\begin{frame}[fragile]
\frametitle{A Coq Computation of Countdown}

\textbf{The worker.}
The \emph{countdown worker}
to $a$ of $f$ is a function $\cdw{f}{a}$~~:
\begin{equation*}
\cdw{f}{a}(n, b) = \begin{cases}
0 & \text{if } b = 0 \vee n\le a \\ 1 +\cdw{f}{a}(f(n), b-1) & \text{if } b \ge 1 \wedge n > a
\end{cases}
\end{equation*}

% Linked by A
\begin{lstlisting}
`\href{https://github.com/inv-ack/inv-ack/blob/7270e64a2600b771f2b1b1b151f7d13fb2ae6c97/countdown.v#L103-L108}{Fixpoint cdn\_wkr}` (a : nat) (f : nat -> nat) (n b : nat) : nat :=
  match b with 0 => 0 | S b' =>
    if (n <=? a) then 0 else S (cdn_wkr f a (f n) k')
end.
\end{lstlisting}

\bigskip

\pause 
\textbf{The Countdown.} Budget $b = n$ is sufficient.
\impinline{Redefine $\cdt{f}{a}(n) \triangleq \cdw{f}{a}(n, n)$.}
% Linked by A
\begin{lstlisting}
`\href{https://github.com/inv-ack/inv-ack/blob/7270e64a2600b771f2b1b1b151f7d13fb2ae6c97/countdown.v#L112}{Definition countdown\_to}` a f n := cdn_wkr a f n n.
\end{lstlisting}
\end{frame}


%\subsection*{}
%\begin{frame}
%\frametitle{Step 2: Inverting the Hierarchies via \textbf{Countdown}}
%\setbeamertemplate{enumerate items}[default]
%\setbeamercolor{local structure}{fg=black}
%
%\begin{enumerate}[\bfseries 1.]
%	\itemsep 3ex
%	\item<done@1->
%	\emph{What is inverse?} \textbf{Upper inverse and increasing functions}.
%	
%	\item<done@1->
%	\emph{Can Repeater preserve Invertibility?} \textbf{Repeatable functions}.
%	
%	\item<done@1->
%	\emph{Computing inverse with inverse:} \textbf{Contractions and Countdown}.
%	
%	\item<come@2->
%	\textbf{Invert each level in hyperoperations/Ackermann hierarchies}.
%\end{enumerate}
%\end{frame}


\subsection{Inverting the Hyperoperations/Ackermann Hierarchies}
\begin{frame}[fragile]
\frametitle{The Inverse Hyperoperation Hierarchy}
The inverse hyperoperations, written $a\angle{n}b$, are defined as:
\begin{equation*}
a\angle{n}b \; \triangleq \; \begin{cases}
b - 1 & \text{ if } n = 0 \\
\cdt{a\angle{n-1}}{a_n}(b) & \text{ if } n \ge 1
\end{cases}
\quad \text{ where } \ a_n = \begin{cases}
a & \hspace{-7pt}\text{ if } n = 1 \\
0 & \hspace{-7pt}\text{ if } n = 2 \\
1 & \hspace{-7pt}\text{ if } n \ge 3
\end{cases}
\end{equation*}
\vspace{-1em}
% linked by A
\pause 
\begin{lstlisting}
`\href{https://github.com/inv-ack/inv-ack/blob/7270e64a2600b771f2b1b1b151f7d13fb2ae6c97/applications.v#L28-L33}{Fixpoint inv\_hyperop}` (a n b : nat) : nat :=
  match n with 0 => b - 1 | S n' =>
    countdown_to (hyperop_init a n') (inv_hyperop a n') b
  end.
\end{lstlisting}

\smallskip

\pause 
Interesting individual levels:
$\href{https://github.com/inv-ack/inv-ack/blob/7270e64a2600b771f2b1b1b151f7d13fb2ae6c97/applications.v#L102-L113}{a\angle{2}b}
= \left\lceil b/a \right\rceil$,
$\href{https://github.com/inv-ack/inv-ack/blob/7270e64a2600b771f2b1b1b151f7d13fb2ae6c97/applications.v#L115-L124}{a\angle{3}b}
= \left\lceil \log_a b \right\rceil$, and
$\href{https://github.com/inv-ack/inv-ack/blob/7270e64a2600b771f2b1b1b151f7d13fb2ae6c97/applications.v#L126-L128}{a\angle{4}b}
= \log^*_a b$ (not currently in Coq standard library).
\end{frame}


\begin{frame}[fragile]
\frametitle{The Inverse Ackermann Hierarchy}

\textbf{Naive approach.} $\Ack_{i+1} = \rf{\big(\Ack_i\big)}{{\color<2->{red} \Ack_i(1)}}\ $.
Thus $\alpha_{i+1} \triangleq \cdt{\big(\alpha_i\big)}{{\color<2->{red} \Ack_i(1)}}\ $?\\ 

\smallskip

\pause
\textcolor{red}{\textbf{Flaw!}} $\alpha_{i+1}$ depends on {\color{red}$\Ack_i$}.

\bigskip

\pause 
\textbf{Observation.} $\Ack_{i+1}(n) = \Ack_i^{(n)}(\Ack_i(1)) = \Ack_i^{(n+1)}(1)$. Thus
\vspace{-0.7em}
\begin{equation*}
\begin{aligned}
\alpha_{i+1}(n) & = \min\left\{m : n\le \Ack_i^{(m+1)}(1)\right\} \pause 
= \min\left\{m : \big( \alpha_i \big)^{(m+1)}(n)\le 1\right\} \pause \\
& = \min\left\{m : \big( \alpha_i \big)^{(m)}\big(\alpha_i(n)\big)\le 1\right\} \pause
= \cdt{\big(\alpha_i\big)}{1}\big(\alpha_i(n)\big)
\end{aligned}
\end{equation*}
\pause 
\textbf{Redefine:}
$
\alpha_i \triangleq
\begin{cases}
\lambda n.(n - 1) & \text{when } i = 0
\\ \big(\cdt{\alpha_{i-1}}{1}\big)\circ \alpha_{i-1} & \text{when } i\ge 1
\end{cases}
$

\smallskip

\pause 
% Linked by A
\begin{lstlisting}
`\href{https://github.com/inv-ack/inv-ack/blob/7270e64a2600b771f2b1b1b151f7d13fb2ae6c97/inv_ack.v#L37-L41}{Fixpoint alpha}` (m x : nat) : nat :=
  match m with 0 => x - 1 | S m' =>
    countdown_to 1 (alpha m') (alpha m' x)
end.
\end{lstlisting}
\end{frame}


%\section{Inverting Hyperoperations and Ackermann}
%\label{sec: inv-hyperop}
%We now use \emph{countdown} to define the inverse
hyperoperation hierarchy, which features elegant new definitions of
division, $\log$, and $\log^{*}$.
We then modify this technique to arrive at the inverse
Ackermann hierarchy.

\subsection{Inverse hyperoperations, including \li{div}, \li{log}, \li{log*}}

\begin{defn} \label{defn: inv-hyperop}
	The inverse hyperoperations, written $a\angle{n}b$, are defined as:
%	\begin{equation}
%	a\angle{n}b \triangleq \begin{cases}
%	b - 1 & \text{if } n = 0 \\
%	\cdt{a\angle{n-1}}{a_n}(b) & \text{if } n \ge 1
%	\end{cases}
%	\ \ \text{ where } \ a_n = \begin{cases}
%	a & \hspace{-10pt}\text{ if } n = 1 \\
%	0 & \hspace{-10pt}\text{ if } n = 2 \\
%	1 & \hspace{-10pt}\text{ if } n \ge 3
%	\end{cases}
%	\end{equation}
% Edited by Linh
	\begin{equation}
  a\angle{n}b \triangleq \begin{cases}
  b - 1 & \hspace{-5pt}\text{if } n = 0 \\
  \cdt{a\angle{n-1}}{a_n}(b) & \hspace{-5pt}\text{if } n \ge 1
  \end{cases}
   \ \text{ where } \ a_n = \begin{cases}
   a & \hspace{-10pt}\text{ if } n = 1 \\
   0 & \hspace{-10pt}\text{ if } n = 2 \\
  1 & \hspace{-10pt}\text{ if } n \ge 3
  \end{cases}
\end{equation}
% linked by A
%\begin{lstlisting}
%`\href{https://github.com/inv-ack/inv-ack/blob/7270e64a2600b771f2b1b1b151f7d13fb2ae6c97/applications.v#L28-L33}{Fixpoint inv\_hyperop}` (a n b : nat) : nat :=
%  match n with 0 => b - 1 | S n' =>
%    countdown_to (hyperop_init a n') (inv_hyperop a n') b
%  end.
%\end{lstlisting}
% Edited by Linh
\begin{lstlisting}
`\href{https://github.com/inv-ack/inv-ack/blob/7270e64a2600b771f2b1b1b151f7d13fb2ae6c97/applications.v#L28-L33}{\color{blue}Fixpoint inv\_hyperop}` (a n b : nat) : nat :=
  match n with 0 => b - 1 | S n' =>
    countdown_to
      (hyperop_init a n') (inv_hyperop a n') b
  end.
\end{lstlisting}
\end{defn}
\hide{
Continuing with the notion of upper-inverses we have used thus far,
we aim for the \emph{ceiling} of division and $\log$ on real numbers, $\left\lceil b/a \right\rceil$ and $\left\lceil \log_ab \right\rceil$. The iterated logarithm is formally defined as the minimum number of times $\log$ needs to be iteratively
applied to the input for the result to become less than or equal to $1$.
Thus it is natural to define it using \emph{countdown}.

\begin{defn} \label{defn: divc} \label{defn: logc} \label{defn: log*}
\[
\divc(a, b) \! \triangleq \! \cdt{\left(\lambda b. (b{-}a)\right)}{0}(b) \quad
\logc(a, b) \! \triangleq \! \cdt{\left(\lambda b. \divc(a, b)\right)}{1}(b) \quad
\log^*(a, b) \! \triangleq \! \cdt{\left(\lambda b. \logc(a, b)\right)}{1}(b)
%\begin{array}{c}
%\divc(a, b) \triangleq \cdt{\left(\lambda b. (b-a)\right)}{0}(b) \qquad
%\logc(a, b) \triangleq \cdt{\left(\lambda b. \divc(a, b)\right)}{1}(b) \\
%\log^*(a, b) \triangleq \cdt{\left(\lambda b. \logc(a, b)\right)}{1}(b)
%\end{array}
\]
\begin{lstlisting}
Definition divc a b := countdown_to 0 (fun n => n - a) b.
Definition logc a b := countdown_to 1 (divc a) b.
Definition logstar a b := countdown_to 1 (logc a) b.
\end{lstlisting}
\end{defn}
}%end hide
\hide{Proving $\forall a, b \ge 1,~\divc(a, b) = \left\lceil b/a \right\rceil$ serves
to prove that all three definitions are correct, since that statement implies
the correctness of  $\logc$, which in turn implies the correctness of $\log^*$. We will prove this in \ref{blah},
together with the correctness of the entire hyperoperation hierarchy.
}
\noindent where the Curried $a\angle{n{-}1}$ denotes the single-variable function
$\lambda b.a\angle{n{-}1}b$.
We now show that $a\angle{n}$ is the inverse \lb to $a[n]$.
First note that $a\angle{0}\in \contract_0$.~Then:
% Linked by Linh
\begin{lem}
	\href{https://github.com/inv-ack/inv-ack/blob/7270e64a2600b771f2b1b1b151f7d13fb2ae6c97/applications.v#L48-L62}{\color{blue}\coq}
$\forall a.~\forall b.~a\angle{1}b = b - a$.
\end{lem}
\begin{proof}
Theorem~\ref{thm: cdt-recursion} applies because $a\angle{0} \in \contract_0 \subseteq \contract_{a}$, giving the intermediate step shown below.
Thereafter we have $a\angle{1}b = b - a$ by induction on $b$.
%\vspace{-0.8em}
%\begin{equation*}
%a\angle{1}b \quad = \quad \cdt{\left(a\angle{0}\right)}{a}(b) \quad = \quad \begin{cases}
%0 & \text{ if } b\le a \\ 1 + a\angle{1}(b - 1) & \text{ if } b\ge a+1 \end{cases} \\
%\end{equation*}
% Edited by Linh
\begin{equation*}
a\angle{1}b \ = \ \cdt{\left(a\angle{0}\right)}{a}(b) \ = \ \begin{cases}
0 & \text{ if } b\le a \\ 1 + a\angle{1}(b - 1) & \text{ if } b\ge a+1 \end{cases} \\
\end{equation*}
\end{proof}
\begin{col} \label{col: inv-hyperop-1-contr1}
$\forall a \ge 1, a\angle{1} \in \contract_1$.
\end{col}
\hide{\begin{col} \label{col: inv-hyperop-234}
{\color{magenta}$\forall a \in \mathbb{N}$}:
 {\color{blue} For all $a \in \mathbb{N}$},
\begin{enumerate}
	\item If $a\ge 1$, $a\angle{2}b = \divc(a, b) \ \forall b$.
	\item If $a\ge 2$, $a\angle{3}b = \logc(a, b) \ \forall b$ and $a\angle{4}b = \logstarc(a, b) \ \forall b$.
\end{enumerate}
\end{col}}%end hide
%Now we are ready to prove the correctness of the inverse hyperoperations.
\noindent \textbf{N.B.} $a\angle{n}b$ is a total function, but it is never actually used for $a = 0$ when $n \ge 2$ or for $a=1$ when $n \ge 3$. For the values
we do care about, we have our inverse:
% Linked by Linh
\begin{thm} \label{thm: inv-hyperop-correct}
	\href{https://github.com/inv-ack/inv-ack/blob/7270e64a2600b771f2b1b1b151f7d13fb2ae6c97/applications.v#L72-L98}{\color{blue}\coq}
When $n\le 1$, or $n \le 2 \wedge a\ge 1$, or $a\ge 2$, then
$a\angle{n} = \left(a[n]\right)^{-1}$.
\end{thm}
\begin{proof}
$\forall n \ge 2,~$ let $a_0 = a, a_1 = 0, a_n = 1$. Define $P$ and $Q$ as:
\begin{equation*}
P(n) \triangleq  \Big(a[n] \in \repeatable_{a_n}\Big)\ \text{ and } \
Q(n) \triangleq  \Big((a\angle{n} = \left(a[n]\right)^{-1} \Big).
\end{equation*}
We have three goals:
\begin{enumerate}
	\item $\forall a.~ Q(0) \wedge Q(1)$ \\ \vspace{-1.2em}
	\item $\forall a \ge 1.~ Q(2)$ \\ \vspace{-1.2em}
	\item $\forall a \ge 2.~ Q(n)$ \\ \vspace{-1.2em}
\end{enumerate}
	Note that \ \mbox{$\forall n.~ a\angle{n+1}$} $= \cdt{a\angle{n}}{a_n}$ ~and \ $a[n+1] = \rf{a[n]}{a_n}$. By Theorem~\ref{thm: cdt-inv-rf},
% Edited by Linh
%\vspace*{-1em}
%\noindent\begin{minipage}{.45\linewidth}
%\begin{equation}
%P(n) \Rightarrow Q(n) \Rightarrow Q(n+1) \label{eq: tmp-induction-1}
%\end{equation}
%\end{minipage}
%\begin{minipage}{.55\linewidth}
%\begin{equation}
%\qquad a_{n} \ge 1 \Rightarrow P(n) \Rightarrow P(n+1) \label{eq: tmp-induction-2}
%\end{equation}
%\end{minipage}
% Edited by Linh
	\begin{align}
	P(n) & \ \Rightarrow \ Q(n) \ \Rightarrow \ Q(n+1) \label{eq: tmp-induction-1} \\
	\qquad a_{n} \ge 1 & \ \Rightarrow \ P(n) \ \Rightarrow \ P(n+1) \label{eq: tmp-induction-2}
	\end{align}
Goal 1: $P(0) \iff \lambda b.(b+1)\in \repeatable_a$ and
$Q(0)\hspace{-0.3em}\iff\hspace{-0.3em}a\angle{0}\hspace{-0.3em}=\hspace{-0.3em}\left(a[0]\right)^{-1} \lb \iff (b-~1\le c \iff b\le c+1)$.
These are both straightforward, and $Q(1)$ holds by~(\ref{eq: tmp-induction-1}). 
Goal 2: we have $a \ge 1$, so
$P(1)$ holds by $P(0)$ and~(\ref{eq: tmp-induction-2}), and
$Q(2)$ holds by $Q(1)$ and~(\ref{eq: tmp-induction-1}).
Goal 3: we have $a\ge 2$, and using~(\ref{eq: tmp-induction-1}) and $Q(0)$
reduces the goal to $P(n)$. Using~(\ref{eq: tmp-induction-2}) and the fact that $\forall n \neq 1.~a_n\ge 1$, the goal reduces to $P(2)$. This unfolds to:
\begin{equation*}
a[2]\in \repeatable_0 \iff \forall b < c.~ab < ac \quad \wedge \quad \forall b \ge 1.~ab \ge b+1\text{,}
\end{equation*}
which is straightforward for $a\ge 2$. Induction on $n$ gives us the third goal.
\end{proof}
%With this theorem, we have proved the correctness of $\divc$, $\logc$ and $\log^*$, as well as the whole inverse hyperoperation hierarchy.
\begin{rem}
Three early hyperoperations are $a[2]b = ab$, $a[3]b = a^b$ and
$a[4]b = \! ^ba$, so, by Theorem~\ref{thm: inv-hyperop-correct}, we can define their inverses $\left\lceil b/a \right\rceil$, $\left\lceil \log_a b \right\rceil$, and $\log^*_a b$ as	
\href{https://github.com/inv-ack/inv-ack/blob/7270e64a2600b771f2b1b1b151f7d13fb2ae6c97/applications.v#L102-L113}{\color{blue}$a\angle{2}b$},	
\href{https://github.com/inv-ack/inv-ack/blob/7270e64a2600b771f2b1b1b151f7d13fb2ae6c97/applications.v#L115-L124}{\color{blue}$a\angle{3}b$}, and	
\href{https://github.com/inv-ack/inv-ack/blob/7270e64a2600b771f2b1b1b151f7d13fb2ae6c97/applications.v#L126-L128}{\color{blue}$a\angle{4}b$}.
Note that the functions $\log_a b$ and $\log^*_a b$ are not in the Coq Standard Library but are one-liners for us:

\begin{lstlisting}
Definition divc a b := inv_hyperop a 2 b.
Definition logc a b := inv_hyperop a 3 b.
Definition logstar a b := inv_hyperop a 4 b.
\end{lstlisting}
\end{rem}



\hide{
\begin{col}
	$\forall a \in \mathbb{N}$:
	\begin{enumerate}
		\item If $a\ge 1$, $\divc(a, b) = \left\lceil \frac{b}{a} \right\rceil$.
		\item If $a\ge 2$, $\logc(a, b) = \left\lceil \log_ab \right\rceil $ and $\logstarc(a, b) = \log^*_ab$.
	\end{enumerate}
\end{col}
\begin{proof}
	The result follows from \cref{col: inv-hyperop-234}, \cref{thm: inv-hyperop-correct} and the first few expansions of the hyperoperation hierarchy: $a[2]b = ab$, $a[3]b = a^b$ and $a[4]b = \ ^ba$.
\end{proof}

\begin{rem}
	\Cref{thm: cdt-recursion} leads to useful recursive formulas for $\divc$, $\logc$ and $\logstarc$:
	\begin{align*}
	\divc(a, b) & = 1 + \divc(a, b - a) & \text{ if } b > 0 &\text{, and } 0 \text{ otherwise}. \\
	\logc(a, b) & = 1 + \logc\left(a, \left\lceil b/a \right\rceil\right) & \text{ if } b > 1 &\text{, and } 0 \text{ otherwise}. \\
	\logstarc(a, b) & = 1 + \divc(a, \left\lceil \log_ab \right\rceil) & \text{ if } b > 1 &\text{, and } 0 \text{ otherwise}.
	\end{align*}
\end{rem}
}

\subsection{The inverse Ackermann hierarchy}
\label{subsec: inv_ack_hier}

% NEW METHOD FOR INVERSE ACKERMANN
Next, we want to use \emph{countdown} to build the \emph{inverse Ackermann hierarchy}, where each
level $\alpha_i$ inverts the level $\Ack_i$.
We know $\Ack_{i+1} = \rf{\Ack_i}{\Ack_i(1)}$\hspace{0.2em}, so the recursive
rule $\alpha_{i+1} \triangleq \cdt{\alpha_i}{\Ack_i(1)}$ is tempting.
But this approach is flawed because it still depends on $\Ack_i$.
Instead, we reexamine the inverse relationship: suppose $\alpha_i = \big(\Ack_i\big)^{-1}$ and $\alpha_{i+1} = \big(\Ack_{i+1}\big)^{-1}$. Then $\Ack_{i+1}(m) = \big(\Ack_{i}\big)^{(m)} \big(\Ack_{i}(1)\big)$. We then have:
\begin{equation} \label{eq: inv-ack-hier-derive}
\alpha_{i+1}(n)\le m \iff n\le \big(\mathcal{A}_i\big)^{(m+1)}(1) \iff \big(\alpha_i\big)^{(m+1)}(n) \le 1
\end{equation}
Equivalently, $\alpha_{i+1}(n) = \min\big\{m : \big( \alpha_i \big)^{(m+1)}(n)\le 1\big\}$, or $\alpha_{i+1}(n) = \cdt{\alpha_i}{1}\big(\alpha_i(n)\big)$. From \eqref{eq: inv-ack-hier-derive} we can thus define the inverse Ackermann hierarchy:
\begin{defn} \label{defn: inv-ack-hier}
	$ \alpha_i \triangleq \begin{cases}
	\lambda n.(n - 1) & \text{if } i = 0
	\\ \big(\cdt{\alpha_{i-1}}{1}\big)\circ \alpha_{i-1} & \text{if } i\ge 1 \end{cases}
$
\end{defn}
Abracadabra! We can express $\alpha$ using the hierarchy \emph{without referring to $\Ack$ itself}:
\begin{thm} \label{thm: inv-ack-new}
	%For all $n, k$, $n\le \Ack(k, k) \!\!\iff \!\! \alpha_k(n)\le k$. Thus
	$\forall n.~ \alpha(n) = \min\big\{k : \alpha_k(n)\le k \big\}$.
\end{thm}


\begin{table}[t]
	\begin{centermath}
		\begin{array}{c | c@{\hspace{1.5em}} c@{\hspace{1.5em}} c@{\hspace{1.5em}} c@{\hspace{1.5em}} c@{\hspace{1.5em}} c@{\hspace{1.5em}} c@{\hspace{1.5em}} c@{\hspace{1.5em}} c@{\hspace{1.5em}} c@{\hspace{1.5em}}}
			          $\diagbox[height=\line]{$\alpha_k$}{$n$}$ & 0\tikzmark{zero} & 1\tikzmark{one} & \tikzmark{twoa}2\tikzmark{two} & \tikzmark{threea}3\tikzmark{three} & \tikzmark{foura}4\tikzmark{four} & \tikzmark{fivea}5\tikzmark{five} & \tikzmark{sixa}6\tikzmark{six} & \tikzmark{sevena}7\tikzmark{seven} & \tikzmark{eighta}8\tikzmark{eight} & \tikzmark{ninea}9\tikzmark{nine} \\
			\hline
			\alpha_0 & 0\tikzmark{zero_succ} & 0\tikzmark{one_succ} & \tikzmark{two_fail}1 & \tikzmark{three_fail}2 & \tikzmark{four_fail1}3 & \tikzmark{five_fail1}4 & \tikzmark{six_fail1}5 & \tikzmark{seven_fail1}6 & \tikzmark{eight_fail1}7 & \tikzmark{nine_fail1}8 \\
			\alpha_1 & 0 & 0 & 0\tikzmark{two_succ} & 1\tikzmark{three_succ} & \tikzmark{four_fail2}2 & \tikzmark{five_fail2}3 & \tikzmark{six_fail2}4 & \tikzmark{seven_fail2}5 & \tikzmark{eight_fail2}6 & \tikzmark{nine_fail2}7 \\
			\alpha_2 & 0 & 0 & 0 & 0 & 1\tikzmark{four_succ} & 1\tikzmark{five_succ} & 2\tikzmark{six_succ} & 2\tikzmark{seven_succ} & \tikzmark{eight_fail3}3 & \tikzmark{nine_fail3}3 \\
			\alpha_3 & 0 & 0 & 0 & 0 & 0 & 0 & 1 & 1 & 1\tikzmark{eight_succ} & 1\tikzmark{nine_succ} \\
		\end{array}
	\end{centermath}
	\caption{Intuition for $\alpha(n)$ defined without $\Ack(n)$.}
	\label{table:inv_intuition}
\end{table}
  \begin{tikzpicture}[overlay, remember picture, yshift=.25\baselineskip, shorten >=.5pt, shorten <=.5pt]
    \draw [->, green] ({pic cs:zero}) to [bend left]({pic cs:zero_succ});
    \draw [->, green] ({pic cs:one}) to [bend left]({pic cs:one_succ});
   	\draw [->, red] ({pic cs:twoa}) to [bend right]({pic cs:two_fail});
   	\draw [->, green] ({pic cs:two}) to [bend left]({pic cs:two_succ});
   	\draw [->, red] ({pic cs:threea}) to [bend right]({pic cs:three_fail});
   	\draw [->, green] ({pic cs:three}) to [bend left]({pic cs:three_succ});
   	\draw [->, red] ({pic cs:foura}) to [bend right]({pic cs:four_fail1});
   	\draw [->, red] ({pic cs:foura}) to [bend right]({pic cs:four_fail2});
   	\draw [->, green] ({pic cs:four}) to [bend left]({pic cs:four_succ});
   	\draw [->, red] ({pic cs:fivea}) to [bend right]({pic cs:five_fail1});
   	\draw [->, red] ({pic cs:fivea}) to [bend right]({pic cs:five_fail2});
   	\draw [->, green] ({pic cs:five}) to [bend left]({pic cs:five_succ});
   	\draw [->, red] ({pic cs:sixa}) to [bend right]({pic cs:six_fail1});
   	\draw [->, red] ({pic cs:sixa}) to [bend right]({pic cs:six_fail2});
   	\draw [->, green] ({pic cs:six}) to [bend left]({pic cs:six_succ});
   	\draw [->, red] ({pic cs:sevena}) to [bend right]({pic cs:seven_fail1});
   	\draw [->, red] ({pic cs:sevena}) to [bend right]({pic cs:seven_fail2});
   	\draw [->, green] ({pic cs:seven}) to [bend left]({pic cs:seven_succ});
   	\draw [->, red] ({pic cs:eighta}) to [bend right]({pic cs:eight_fail1});
   	\draw [->, red] ({pic cs:eighta}) to [bend right]({pic cs:eight_fail2});
   	\draw [->, red] ({pic cs:eighta}) to [bend right]({pic cs:eight_fail3});
   	\draw [->, green] ({pic cs:eight}) to [bend left]({pic cs:eight_succ});
   	\draw [->, red] ({pic cs:ninea}) to [bend right]({pic cs:nine_fail1});
   	\draw [->, red] ({pic cs:ninea}) to [bend right]({pic cs:nine_fail2});
   	\draw [->, red] ({pic cs:ninea}) to [bend right]({pic cs:nine_fail3});
   	\draw [->, green] ({pic cs:nine}) to [bend left]({pic cs:nine_succ});
  \end{tikzpicture}

% \vspace{-3em}
Table~\ref{table:inv_intuition} gives intuition for our new definition of $\alpha$.
This table is populated exactly as suggested by 
Table~\ref{table: hyperop-ack-inv}:
$\alpha_{0}(n)$ is $n-1$, 
$\alpha_{1}(n)$ is $n-2$,
$\alpha_{2}(n)$ is $\left\lceil \frac{n-3}{2} \right\rceil$, and
$\alpha_{3}(n)$ is $\left\lceil \log_2 ~ (b + 3)\right\rceil - 3$.
Now for any $n$, we search down its column, which is indexed by $k$, 
looking for the smallest $k$ such that $\alpha_k(n) \le k$.

Consider $n = 8$. $\alpha_{0}(8) = 7$. Because $7 \not\le 0$, $\alpha_0$ is rejected.
Similarly $\alpha_{1}(8) = 6 \not\le 1$ and $\alpha_{2}(8) = 3 \not\le 2$
are rejected. Finally $\alpha_{3}(8) = 1$ is accepted because $1 \le 3$. 
Indeed, $\alpha(8) = 3$.
Similar reasoning applies to the other values for $n$. To aid intuition in this table,
we show show unsuccessful searches in {\color{red}red} and successful searches in 
{\color{green}green}.

Now all that remains is to provide a structurally-recursive function that computes $\alpha$.
%\begin{defn} \label{defn: inv-ack-worker}
%	The inverse Ackermann worker is a function $\alpha^{\W}$: %(\mathbb{N}\to \mathbb{N}) \times \mathbb{N}^3\to \mathbb{N}$ such that for all $n, k, b\in \mathbb{N}$ and $f:\mathbb{N}\to \mathbb{N}$:
%	\begin{equation} \label{eq: inv-ack-worker-recursion}
%	\alpha^{\W}(f, n, k, b) \triangleq \begin{cases}
%	k & \text{if } b = 0 \vee n\le k \\ \alpha^{\W}(\cdt{f}{1}\circ f , \cdt{f}{1}(n), k+1, b-1) & \text{if } b \ge 1 \wedge n \ge k+1
%	\end{cases}
%	\end{equation}
%\end{defn}
% Edited by Linh
\begin{defn} \label{defn: inv-ack-worker}
	The inverse Ackermann worker, written $\alpha^{\W}$, is a function from $\mathbb{N}^4$ to $\mathbb{N}$ defined as: %(\mathbb{N}\to \mathbb{N}) \times \mathbb{N}^3\to \mathbb{N}$ such that for all $n, k, b\in \mathbb{N}$ and $f:\mathbb{N}\to \mathbb{N}$:
	\begin{equation} \label{eq: inv-ack-worker-recursion}
	\begin{aligned}
	& \alpha^{\W}(f, n, k, b) \\ 
	& \triangleq \begin{cases}
	k & \text{if } b = 0 \vee n\le k \\ \alpha^{\W}(\cdt{f}{1}\circ f , \cdt{f}{1}(n), k+1, b-1) & \text{if } b \ge 1 \wedge n \ge k+1
	\end{cases}
		\end{aligned}
	\end{equation}
\end{defn}
\noindent Next, we show that this function computes the inverse Ackermann function when passed appropriate arguments.
\hide{
Given the arguments $\big(\alpha_i, \alpha_i(n), i, b - i\big)$ such that $\alpha_i(n) > i$ and $b > i$, $\alpha^{\W}$ takes on arguments $\big(\alpha_{i+1}, \alpha_{i+1}(n), i+1, b - (i+1)\big)$ at the next recursive call. Thus if $\alpha^{\W}$ is given a sufficient budget $b$, it will recursively transform the tuple $(\alpha_k, \alpha_k(n), k, b - k)$ until a point $k$ where $\alpha_k(n)\le k$, and will then return $k$. We now need to show that $\alpha^{\W}$ correctly computes $\alpha(n)$ given a reasonable budget.
The following theorem demonstrates that a budget of $n$ suffices.
}%end hide
\hide{
\noindent We are finally ready for a strategy to compute the inverse Ackermann function:
}%end hide
%formalizes a setting for $\W\alpha$ to work.
% Linked by Linh
\begin{thm} \label{thm: inv-ack-worker-correct}
	\href{https://github.com/inv-ack/inv-ack/blob/7270e64a2600b771f2b1b1b151f7d13fb2ae6c97/inv_ack.v#L199-L231}{\color{blue}\coq}
	$\forall n.~\alpha^{\W}\big(\alpha_0, \alpha_0(n), 0, n\big) = \alpha(n)$.
\end{thm}
\begin{proof}[Proof outline]
	When given the arguments $\big(\alpha_i, \alpha_i(n), i, b - i\big)$ such that $\alpha_i(n) > i$ and $b > i$, $\alpha^{\W}$ takes on the arguments $\big(\alpha_{i+1}, \alpha_{i+1}(n), i+1, b - (i+1)\big)$ at the next recursive call. A simple induction on $k$ then shows that if $k\le \min\big\{b, \alpha_k(n)\big\}$,
	\begin{equation} \label{eq: inv-ack-worker-intermediate}
	\alpha^{\W}\big(\alpha_0, \alpha_0(n), 0, b\big) = \alpha^{\W}\big(\alpha_k, \alpha_k(n), k, b-k\big)
	\end{equation}
	Let $m \triangleq \min\big\{k : \alpha_k(n) \le k \}$. Then $m\le n$ since $\alpha_n(n)\le n$. \eqref{eq: inv-ack-worker-intermediate} then implies:
	$$ \alpha^{\W}\big(\alpha_0, \alpha_0(n), 0, n\big) = \alpha^{\W}\big(\alpha_m, \alpha_m(n), m, n - m\big) = m = \alpha(n) $$
\end{proof}
\noindent We put a more involved mathematical proof of correctness in Appendix~\ref{apx:proof_correct_inv_ack_worker}, and a mechanized proof
% Linked by Linh
	\href{https://github.com/inv-ack/inv-ack/blob/7270e64a2600b771f2b1b1b151f7d13fb2ae6c97/inv_ack.v#L163-L231}{\color{blue}here}.
We thus have a (re-)definition of inverse Ackermann that is %computationally
definable via a Coq-accepted worker, \emph{i.e.} $\alpha(n) \triangleq \alpha^{\W}\big(\alpha_0, \alpha_0(n), 0, n\big)$. Below, we present the inverse
Ackermann function in Gallina.

% Linked by A
%\begin{lstlisting}
%`\href{https://github.com/inv-ack/inv-ack/blob/7270e64a2600b771f2b1b1b151f7d13fb2ae6c97/inv_ack.v#L155-L161} {Fixpoint inv\_ack\_wkr}` (f : nat -> nat) (n k b : nat) : nat :=
%  match b with 0 => 0 | S b' =>
%    if (n <=? k) then k else let g := (countdown_to f 1) in
%      inv_ack_wkr (compose g f) (g n) (S k) b
%  end.
%
%`\href{https://github.com/inv-ack/inv-ack/blob/7270e64a2600b771f2b1b1b151f7d13fb2ae6c97/inv_ack.v#L37-L41}{Fixpoint alpha}` (m x : nat) : nat :=
%  match m with 0 => x - 1 | S m' =>
%    countdown_to 1 (alpha m') (alpha m' x)
%  end.
%
%`\href{https://github.com/inv-ack/inv-ack/blob/7270e64a2600b771f2b1b1b151f7d13fb2ae6c97/inv_ack.v#L167}{Definition inv\_ack}` := inv_ack_wkr (alpha 0) (alpha 0 n) 0 n.
%\end{lstlisting}
% Edited by Linh
\begin{lstlisting}
`\href{https://github.com/inv-ack/inv-ack/blob/7270e64a2600b771f2b1b1b151f7d13fb2ae6c97/inv_ack.v#L155-L161} {\color{blue}Fixpoint inv\_ack\_wkr}` f n k b :=
  match b with 0 => 0 | S b' =>
    if (n <=? k) then k
      else let g := (countdown_to f 1) in
        inv_ack_wkr (compose g f) (g n) (S k) b
  end.

`\href{https://github.com/inv-ack/inv-ack/blob/7270e64a2600b771f2b1b1b151f7d13fb2ae6c97/inv_ack.v#L37-L41}{\color{blue}Fixpoint alpha}` m x :=
  match m with 0 => x - 1 | S m' =>
    countdown_to 1 (alpha m') (alpha m' x)
  end.

`\href{https://github.com/inv-ack/inv-ack/blob/7270e64a2600b771f2b1b1b151f7d13fb2ae6c97/inv_ack.v#L167}{\color{blue}Definition inv\_ack}` :=
  inv_ack_wkr (alpha 0) (alpha 0 n) 0 n.
\end{lstlisting}

\noindent Note that this is not the linear-time computation we presented in
Figure~\ref{fig:standalone}. We will arrive at that code via an improvement discussed in the next section.
% OLD METHOD FOR INVERSE ACKERMANN

%Using the Ackermann kludge from \cref{sec: overview}, we sketch a method to find the inverse Ackermann function from the inverse hyperoperations.
%\begin{thm} \label{thm: inv-hyperop-inv-ack}
%	For all $n, k$, $n\le \Ack(k, k) \iff 2\angle{k}(n+3)\le k+3$.
%\end{thm}
%\begin{proof}
%	$n\le \Ack(k, k) \iff n\le 2[k](k+3) - 3$. Now $2[k]$ is an expansion by \cref{thm: inv-hyperop-correct}'s proof, so $2[k](k+3) \ge k+3\ge 3$, so  $n\le 2[k](k+3) - 3 \iff n+3 \le 2[k](k+3)\iff 2\angle{k}(n+3)\le k+3$, again by \cref{thm: inv-hyperop-correct}.
%\end{proof}
%\Cref{thm: inv-hyperop-inv-ack} allows a simple method to compute the Inverse Ackermann function, based on the \emph{budget} idea in \emph{countdown worker}.
%\begin{defn} \label{defn: inv-ack-worker}
%	The \emph{inverse Ackermann worker} of $f$ is a function $\W\alpha\ : \mathbb{N}^4\to \mathbb{N}$ such that for all $n, k, b\in \mathbb{N}$ and $f:\mathbb{N}\to \mathbb{N}$:
%	$$ \W\alpha(f, n, k, b) = \begin{cases}
%	k - 3 & \text{if } b = 0 \vee f(n)\le k+3 \\ \W\alpha(\cdt{f}{a_k}\ , n, k+1, b-1) & \text{if } b \ge 1 \wedge f(n) > k+3
%	\end{cases} $$
%	where $a_0 = 2$, $a_1 = 0$ and $a_k = 1 \ \forall k\ge 2$.
%\end{defn}
%The function $\W\alpha$ takes a function $f$ and pretends it was the first level of the inverse hyperoperation hierarchy. It then keeps applying countdown to $f$, supposedly arrives at the $\text{m}^{\text{th}}$-level at the $\text{m}^{\text{th}}$-recursive step, until the budget $b$ is exhausted or $f(n)\le k+3$, i.e. \emph{early stopping}, and will output $k$. If at the beginning $k=3$ and $f = 2\angle{0}$, the early stopping condition becomes $2\angle{k}n \le k+3$, which when replace $n$ by $n+3$ gives what we need in \cref{thm: inv-hyperop-inv-ack}.
%The inverse Ackermann function can be defined as follows.
%\begin{thm} For all $n$, we have $\alpha(n) = \W\alpha(\lambda m.(m - 1), n+3, 3, n)$.
%\end{thm} 

\section*{Roadmap}
\begin{frame}
\frametitle{Roadmap}
\setlength{\leftmargini}{3.8em}
\setbeamertemplate{enumerate items}[default]
\setbeamercolor{local structure}{fg=black}

\imppar{\textbf{Goal.} Inverting $\Ack$ - without computing $\Ack$.}

\bigskip

\begin{enumerate}[\bfseries Step 1.]
	\itemsep 2ex
	\item<done@1->
	Build the hyperoperations/Ackermann hierarchies via \textbf{Repeater}.
	
	\item<done@1->
	Build inverses of the hyperoperations/Ackermann hierarchies via \textbf{Countdown}, the inverse of \emph{Repeater}.
	
	\item<come@2->
	Use the hierarchy to implement the inverse Ackermann function \\ for inputs encoded in unary. $\bigO (n)$.
\end{enumerate}
\end{frame}


\section{Step 3: The Inverse Ackermann function in linear time for unary encoding system}
\label{sec: inv-ack}
\subsection{Using the Hierarchy to Implement $\alpha(n)$}
\begin{frame}
\frametitle{Implementation Idea}
\textbf{Redefinition.}
$
\alpha(n) = \min\{k: n\le A(k, k) \} \triangleq \min\{k: \alpha_k(n)\le k \}
$

\bigskip

\pause 
\textbf{The worker function.} $\begin{cases}
\text{Budget } b: & \hspace{-1em}\text{Maximum } b \text{ steps. } \\
\text{Step } i: & \hspace{-1em}\text{Compute } \alpha_i(n). \\
\text{Stops when: } & \hspace{-1em}\text{budget reaches } 0 \text{ or } \alpha_i(n) \le i.
\end{cases}$

%\textbf{The worker function.}
%\bigskip
%\begin{itemize}
%	\item Budget $b$: Maximum $b$ steps.
%	\item Step $i$: Computes $\alpha_i(n)$.
%	\item Stops when budget reaches $0$ or $\alpha_i(n) \le i$.
%\end{itemize}

\bigskip

\pause 
\textbf{Advantage.} $\alpha_{i+1}(n) = \cdt{\big(\alpha_i\big)}{1}\big(\alpha_i(n)\big)$\\
\smallskip
So the next step is computable from the previous.
\end{frame}


\begin{frame}[fragile]
\frametitle{The Worker Function}

\uncover<1->{$\displaystyle 	\alpha^{\W}(f, n, k, b) = \begin{cases}
k & \text{if } b = 0 \vee n\le k \\ \alpha^{\W}(\cdt{f}{1}\circ f , \cdt{f}{1}(n), k+1, b-1) & \text{if } b \ge 1 \wedge n \ge k+1
\end{cases}$}

\bigskip

\begin{overlayarea}{\linewidth}{3cm}
\begin{onlyenv}<2->
	\begin{lstlisting}
`\href{https://github.com/inv-ack/inv-ack/blob/7270e64a2600b771f2b1b1b151f7d13fb2ae6c97/inv_ack.v#L155-L161} {Fixpoint inv\_ack\_wkr}` (f : nat -> nat) (n k b : nat) : nat :=
match b with 0 => 0 | S b' =>
  if (n <=? k) then k else let g := (countdown_to f 1) in
                      inv_ack_wkr (compose g f) (g n) (S k) b
end.
\end{lstlisting}
\end{onlyenv}
\end{overlayarea}

\end{frame}



\begin{frame}
\frametitle{The Worker Function}

\textbf{Observation.}
\begin{equation*}
\begin{array}{lllllll}
Initial \ Arguments & & \big(\alpha_i, & \alpha_i(n), & i, & b - i \big) \\[5pt]
\pause
b - i > 0, \ \alpha_i(n) > i & \pause \to & \big(\cdt{\alpha_i}{1}\circ \alpha_i, & \cdt{\alpha_i}{1}(\alpha_i(n)), & i + 1, & b - i - 1\big) \\[5pt]
\pause
{\color{blue} b > i}, \ \hspace{15pt} \alpha_i(n) > i & \to  & \big({\color{blue} \alpha_{i+1}}, & {\color{blue} \alpha_{i+1}(n)}, & i + 1, & {\color{blue} b - (i + 1)}\big) \\
\end{array}
\end{equation*}

\pause
\textbf{Execution.}
\begin{equation*}
\begin{array}{l|lllllll}
Step & Initial \ Arguments &  &  \big(\alpha_0, & \alpha_0(n), & 0, & b - 0 \big) \\[3pt]
\pause
0 & b > 0, \ \alpha_0(n) > 0 & \pause \to & \big(\alpha_1, & \alpha_1(n), & 1, & b - 1\big) \\[3pt]
\pause
1 & b > 1, \ \alpha_1(n) > 1 & \pause \to  & \big(\alpha_2, & \alpha_2(n), & 2, & b - 2\big) \\[3pt]
\pause
\vdots & \vdots & \vdots & & \multicolumn{3}{l}{\vdots} \\
%b > k-1, \ \alpha_{k-1}(n) > k - 1 & \to  & \big(\alpha_k, & \alpha_k(n), & k, & b - k\big) & k-1 \\
\pause
k & b > k, \ {\color{red} \alpha_{k}(n) \le k} & \pause \to  & \multicolumn{4}{l}{\impinline{$k = \alpha(n)$}}
\end{array}
\end{equation*}
\end{frame}



\begin{frame}
\frametitle{The Inverse Ackermann Function}

Any budget $b \ge \alpha(n)$ suffices, so we can choose $b = n$.

\smallskip

\pause
\textbf{First version:}.
\pause
\begin{equation*}
\alpha(n) = \alpha^{\mathcal{W}}\big(\alpha_0, \alpha_0(n), 0, n\big)
= \alpha^{\mathcal{W}}\big(\lambda n(n - 1), n - 1, 0, n\big)
\end{equation*}

\pause
\textbf{Time Complexity.}

\vspace{5pt}

\begin{tabular}{l|l|l}
Computation & Time & Reason \\ \hline
$\topspace{2pt} \alpha_0(n) = n - 1$ & $\Theta(1)$ & Coq SL \\
$\topspace{2pt} \alpha_1(n) = n - 2$ & $\Theta(n)$ & Apply $\alpha_0$ for $n - 1$ times \\
$\topspace{2pt} \alpha_2(n) = \left \lceil \frac{n-3}{2} \right \rceil$ & $\Theta(n^2)$ & Apply $\alpha_1$ for $n/2$ times
\end{tabular}

\vspace{7pt}

\textcolor{red}{\textbf{Total time is at least $\mathbf{\Omega(n^2)}$!}}

\end{frame}


\begin{frame}
\frametitle{The Inverse Ackermann Function}

\textbf{Simple improvement.}
\pause
\begin{equation*}
\begin{aligned}
\alpha(n) & = \begin{cases}
\alpha^{\mathcal{W}}\big(\alpha_{{\color{red} 1}}, \alpha_{{\color{red} 1}}(n), {{\color{red} 1}}, n - {{\color{red} 1}}\big) & \hspace{3em}\text{ when } n > \Ack(0) \\
0 & \hspace{3em}\text{ when } n \le \Ack(0)
\end{cases} \\[5pt]
\pause \alpha(n) & = \begin{cases}
\alpha^{\mathcal{W}}\big(\lambda n(n - 2), n - 2, 1, n - 1\big) & \text{ when } n > 1 \\
0 & \text{ when } n \le 1
\end{cases}
\end{aligned}
\end{equation*}

\pause
Now, computing $\alpha_1(n) = n - 2$ takes time $\Theta(1)$,\\
so computing $\alpha_2(n)$ takes time $\Theta(n)$.

\bigskip
\pause
\textcolor{red}{\textbf{In fact, the total time is still $\mathbf{\Theta(n)}$.}}

\end{frame}


%\begin{frame}
%\frametitle{Time Complexity of $\alpha$: A Sketch}
%
%$
%{\color{green}\alpha_{i+1}} = \cdt{\big({\color{blue}\alpha_{i}}\big)}{1}\circ {\color{blue}\alpha_i} \quad \text{ and } \quad {\color{green} \alpha_{i+1}(n)} = \cdt{\big({\color{blue}\alpha_{i}}\big)}{1}\big({\color{red}\alpha_{i}(n)}\big).
%$
%
%\bigskip
%
%$ \displaystyle
%\begin{aligned}
%\alpha^{\mathcal{W}}\big({\color{blue}\alpha_i}, \ {\color{red}\alpha_i(n)}, \ i, \ b\big)
%& = \alpha^{\mathcal{W}}\left(
%\cdt{{\color{blue}\alpha_i}}{1}\circ {\color{blue}\alpha_i},  \
%\cdt{{\color{blue}\alpha_i}}{1}\big({\color{red}\alpha_i(n)}\big), \
%i+1, \ b - 1\right) \\
%& = \alpha^{\mathcal{W}}\left(
%{\color{green}\alpha_{i+1}}, \ \ \ \
%{\color{green}\alpha_{i+1}(n)}, \ \ \ \ \
%i+1, \ b - 1\right).
%\end{aligned}
%$
%
%\bigskip
%
%The recursive formula of $\alpha^{\mathcal{W}}$ contains the recursive formula of $\alpha_i$.
%
%%The computation $\alpha(n)$ up to step $i$ incorporates the computation of $\alpha_i(n)$, plus some minor computations.
%
%\bigskip
%
%Thus, most of $\alpha(n)$'s running time is to compute $\alpha_k(n)$, where $k \triangleq \alpha(n)$.
%
%\bigskip
%
%Similarly, for each $i$, most of $\alpha_i(n)$'s running time is to compute the chain $\alpha_i(n), \alpha_i^{(2)}(n), \alpha_i^{(3)}(n) \ldots, $ until reaching $1$.
%
%%\begin{equation*}
%%\begin{array}{l|lllllll}
%%Step & Initial \ Arguments &  &  \big(\alpha_0, & \alpha_0(n), & 0, & b - 0 \big) \\[3pt]
%%0 & b > 0, \ \alpha_0(n) > 0 & \to & \big(\alpha_1, & \alpha_1(n), & 1, & b - 1\big) \\[3pt]
%%1 & b > 1, \ \alpha_1(n) > 1 & \to  & \big(\alpha_2, & \alpha_2(n), & 2, & b - 2\big) \\[3pt]
%%\vdots & \vdots & \vdots & & \multicolumn{3}{l}{\vdots} \\
%%%b > k-1, \ \alpha_{k-1}(n) > k - 1 & \to  & \big(\alpha_k, & \alpha_k(n), & k, & b - k\big) & k-1 \\
%%k & b > k, \ \alpha_{k}(n) \le k & \to  & \multicolumn{4}{l}{k = \alpha(n)}
%%\end{array}
%%\end{equation*}
%
%\end{frame}



\begin{frame}
\frametitle{Time Complexity of $\alpha$: A Sketch}

Let $\runtime_f(n)$ denotes the running time of computing $f(n)$ given $f$ and $n$.
\pause
\begin{equation*}
\runtime_{\alpha_{i+1}}(n) \approx
\runtime_{\alpha_i}(n) + \runtime_{\alpha_i}\big(\alpha_i(n)\big)
+ \runtime_{\alpha_i}\big(\alpha_i^{(2)}(n)\big) + \ldots (\text{until 1})
\end{equation*}

\smallskip

\pause
Manual computation: $\runtime_{\alpha_3}(n)\le 4n + 4$.

\smallskip

\pause
Suppose $\runtime_{\alpha_i}(n) \le 4n + \bigO\big(\log_2n\big)$. \pause Then
\begin{equation*}
\begin{aligned}
& \runtime_{\alpha_{i+1}}(n) \ \lessapprox \ 4n + \bigO\big(\log_2n\big) + 4\log_2n + \bigO\big(\log_2^{(2)}n\big) + \ldots \\
& \quad = 4n + \bigO\left(\log_2n + \log_2^{(2)}n + \log_2^{(3)}n + \ldots\right)
 = 4n + \bigO\big(\log_2n\big)
\end{aligned}
\end{equation*}
\pause
Therefore, \imppar{$T_{\alpha}(n) \approx \runtime_{\alpha_{\alpha(n)}}(n) \lessapprox 4n + \bigO\big(\log_2n\big) \lessapprox \Theta(n)$} by induction.

\end{frame}



\section*{Roadmap}
\begin{frame}
\frametitle{Roadmap}
\setlength{\leftmargini}{3.8em}
\setbeamertemplate{enumerate items}[default]
\setbeamercolor{local structure}{fg=black}

\imppar{\textbf{Goal.} Inverting $\Ack$ - without computing $\Ack$.}

\bigskip

\begin{enumerate}[\bfseries Step 1.]
	\itemsep 2ex
	\item<done@1->
	Build the hyperoperations/Ackermann hierarchies via \textbf{Repeater}.
	
	\item<done@1->
	Build inverses of the hyperoperations/Ackermann hierarchies via \textbf{Countdown}, the inverse of \emph{Repeater}.
	
	\item<done@1->
	Use the hierarchy to implement the inverse Ackermann function \\ for inputs encoded in unary. $\bigO (n)$.
\end{enumerate}
\end{frame}


\section*{Roadmap}
\begin{frame}
\frametitle{Roadmap}
\setlength{\leftmargini}{3.8em}
\setbeamertemplate{enumerate items}[default]
\setbeamercolor{local structure}{fg=black}

\imppar{\textbf{Goal.} Inverting $\Ack$ - without computing $\Ack$.}

\bigskip

\begin{enumerate}[\bfseries Step 1.]
	\itemsep 2ex
	\item<done@1->
	Build the hyperoperations/Ackermann hierarchies via \textbf{Repeater}.
	
	\item<done@1->
	Build inverses of the hyperoperations/Ackermann hierarchies via \textbf{Countdown}, the inverse of \emph{Repeater}.
	
	\item<done@1->
	Use the hierarchy to implement the inverse Ackermann function \\ for inputs encoded in unary. $\bigO (n)$.
\end{enumerate}
\end{frame}

%\section{Further Discussion}
%\label{sec: discussion}
%\newcommand{\ackt}{\ensuremath{\hat{\alpha}}}

\subsection{The value of a linear-time solution to the hierarchy}

Our functions' linear runtimes can be understood in two distinct but
complementary ways.  A runtime less than the bitlength is impossible
without prior knowledge of the size of the input.  Accordingly, in
an information-theory or pure-mathematical sense, our definitions are
optimal up to constant factors.  And of course in practice, linear-time
solutions are highly usable in real computations.

Sublinear solutions are possible with \emph{a priori} knowledge about
the function and bounds on the inputs one will receive.
An extreme case is $\alpha(n)$, which has value $4$ for all practical
inputs greater than $61$. Accordingly,
this function can be inverted in $O(1)$ in practice.  That said, 
such solutions require external knowledge of the problems and
lookup tables within the algorithm to store precomputed
values, and thus fall more into the realm of engineering than mathematics. 

\subsection{The two-parameter inverse Ackermann function}
Some authors~\cite{chazelle,tarjan} prefer a two-parameter inverse Ackermann function.
\begin{defn} \label{defn: 2para-alpha}
	The two-parameter inverse Ackermann function is defined as:
	\begin{equation} \label{eq: tmp-2para-alpha}
	\ackt (m, n) \triangleq \min\left\{i \ge 1 : \Ack\left(i, \left\lfloor \frac{m}{n} \right\rfloor \right)\ge \log_2n \right\}
	\end{equation}
\end{defn}
Note that $\ackt(n, n)$ and the single-parameter $\alpha(n)$
are neither equal nor directly related, but
it is straightforward to modify our techniques to compute $\ackt(m, n)$.
\hide{This function arises from deep runtime analysis of the disjoint-set data structure. Tarjan \cite{tarjan} showed that, in the disjoint-set data structure, the time required $t(m,n)$ for a sequence of $m$ \textsc{\color{magenta}FIND}s intermixed with $n-1$ \textsc{\color{magenta}UNION}s (such that $m \geq n$) is bounded as: $k_{1}m\cdot\alpha(m,n) \leq t(m,n) \leq k_{2}m\cdot\alpha(m,n)$. In graph theory, Chazelle \cite{chazelle} showed that the minimum spanning tree of a connected graph with $n$ vertices and $m$ edges can be found in time $O(m\cdot\alpha(m,n))$. Computing this function is in fact easier than $\alpha(n)$, as when $m$ and $n$ are given, we are reduced to finding the minimum $i\ge 1$ such that $\Ack_i(s)\ge t$ for $s, t$ fixed, which can be done with the following variant of the \emph{inverse Ackermann worker}.
}
\begin{defn} \label{defn: 2para-inv-ack-worker}
	The {two-parameter inverse Ackermann worker}
	,written $\ackt^{\W}$, is a function $\mathbb{N}^4\to \mathbb{N}$, defined by:
	\hide{$(\mathbb{N}\to \mathbb{N}) \times \mathbb{N}^3\to \mathbb{N}$ such that for all $n, k, b\in \mathbb{N}$ and $f:\mathbb{N}\to \mathbb{N}$:}
%	\begin{equation} \label{eq: 2para-inv-ack-worker-recursion}
%	\ackt^{\W}(f, n, k, b) = \begin{cases}
%	0 & \text{if } b = 0 \vee n\le k \\ 1 + \ackt^{\W}\big(\cdt{f}{1}\circ f , \cdt{f}{1}(n), k, b-1\big) & \text{if } b \ge 1 \wedge n \ge k+1
%	\end{cases}
%	\end{equation}
  \begin{equation} \label{eq: 2para-inv-ack-worker-recursion}
  \begin{aligned}
  & \ackt^{\W}(f, n, k, b) \\
  & \triangleq \begin{cases}
  0 & \text{if } b = 0 \vee n\le k \\ 1 + \ackt^{\W}\big(\cdt{f}{1}\circ f , \cdt{f}{1}(n), k, b-1\big) & \text{if } b \ge 1 \wedge n \ge k+1
  \end{cases}
  \end{aligned}
  \end{equation}
\end{defn}
%Similar to the one-parameter version, the following theorem establishes the correct setting for $\W\alpha_2$ to compute $\alpha(m, n)$.
%\begin{thm}
%	$\displaystyle \ackt(m, n) = 1 + \ackt^{\W}\left(\alpha_1, \alpha_1\big(\lceil\log_2n \rceil\big), \left\lfloor \frac{m}{n} \right\rfloor, \lceil\log_2n \rceil \right)$.
%\end{thm}
% Edited by Linh
\begin{thm} For all $m$ and $n$,
	\begin{equation*}
	\displaystyle \ackt(m, n) = 1 + \ackt^{\W}\left(\alpha_1, \alpha_1\big(\lceil\log_2n \rceil\big), \left\lfloor \frac{m}{n} \right\rfloor, \lceil\log_2n \rceil \right).
	\end{equation*}
\end{thm}
We mechanize the above for both \href{https://github.com/inv-ack/inv-ack/blob/7270e64a2600b771f2b1b1b151f7d13fb2ae6c97/inv_ack.v#L245-L248}{\color{blue}unary} and \href{https://github.com/inv-ack/inv-ack/blob/7270e64a2600b771f2b1b1b151f7d13fb2ae6c97/bin_inv_ack.v#L222-L228}{\color{blue}binary} inputs in our codebase.

\hide{
	\begin{proof}[Proof Sketch]
		It is easy to prove in a similar fashion to \cref{lem: inv-ack-worker-intermediate} that for all $n, b, k$ and $i$, if $\alpha_i(n) > k$ and $b > i$, then
		\begin{equation*}
		\W\alpha_2\big(\alpha_1, \alpha_1(b), k, b\big) = i + \W\alpha_2\big(\alpha_{i+1}, \alpha_{i+1}(n), k, b - i\big)
		\end{equation*}
		Now let $k \triangleq \lfloor m/n \rfloor$, $b \triangleq \lceil \log_2n \rceil$ and $l \triangleq \min\big\{i : \alpha_i(b)\le k\big\}$, which exists because $\Ack(i, \cdot)$ increases strictly with $i$. Then $\alpha(m, n) = \max{1, l}$. If $l = 0$, $\alpha_1(b) \le \alpha_0(b) \le k$, so $\W\alpha_2\big(\alpha_1, \alpha_1(b), k, b\big) = 0$, as desired. If $l \ge 1$,
		\begin{equation*}
		1 + \W\alpha_2\big(\alpha_1, \alpha_1(b), k, b\big)
		= 1 + l - 1 + \W\alpha_2\big(\alpha_l, \alpha_l(b), k, b-l\big) = l
		\end{equation*}
		Here $b\ge l$ due to the fact that $\Ack(b, k)\ge b$, so $\alpha_b(b)\le k$. This completes the proof.
\end{proof}}%end hide




\section*{Conclusion}
\begin{frame}
\frametitle{Conclusion}
%\frametitle{Abstract}
%TODO: Past tense
We inverted the hyperoperation 
and Ackermann hierarchies \\ and implemented 
our inverse hierarchies in Gallina.
We used our inverses to compute the upper inverse of the 
diagonal Ackermann function. 

\bigskip

Our computations run in linear time
-- $\Theta(b)$, where $b$=bitlength -- \\ for inputs represented in 
both unary and binary.

\bigskip

We showed that these functions are consistent with the usual \\ definition of the inverse Ackermann function~$\alpha(n)$.
\keywords{Ackermann  \and hyperoperations \and Coq}
\end{frame}

\section*{Demo}
\begin{frame}
\frametitle{Demo}
We present a brief demonstration of our time bounds in Coq.
\end{frame}

\newcounter{totalframenumber}
\setcounter{totalframenumber}{\inserttotalframenumber}

\appendix


\addtocounter{framenumber}{\value{totalframenumber}}

\setbeamertemplate{footline}{%
	\hspace{1em}\small\arabic{framenumber}/\arabic{totalframenumber}
	\hfill\usebeamertemplate***{navigation symbols}}
	


\section{The Inverse Ackermann function in logarithmic time for binary encoding system}
\label{sec: binary}

\subsection{Overview}

%\begin{frame}
%\frametitle{Review: Binary Encoding System}
%
%\begin{tabular}{l|c|c}
%	Attribute & Unary & Binary Encoding \\ \hline
%	Type in Coq & \texttt{nat} & \texttt{N} \topspace{1pt} \\[6pt]
%	Number $n$ & $\Theta(n)$ bits & $\Theta\big(\log_2n\big)$ bits \\[7pt]
%	Compare $x, y$ & $\Theta\big(\min\{x, y\}\big)$ time & $\Theta\big(\min\{ \log_2x, \log_2y\}\big)$ time \\[9pt]
%	Successor & \multirow{2}{*}{$\Theta(1)$ time} & Worse case: $\Theta\big(\log_2x\big)$ time \\[1pt]
%	function  &             & Amortized: $\Theta(1)$ time \\
%	                    
%\end{tabular}
%
%\bigskip
%
%Binary encoding provides faster computations for many operations.
%
%\smallskip
%
%\pause 
%What about inverse Ackermann? \imppar{\textbf{Goal:} $\bigO \big(\log_2n\big)$}.
%
%\end{frame}


\subsection{Countdown in Binary}


%\begin{frame}[fragile]
%\frametitle{Translating Countdown}
%
%\textbf{The worker.}
%\begin{lstlisting}
%`\href{https://github.com/inv-ack/inv-ack/blob/7270e64a2600b771f2b1b1b151f7d13fb2ae6c97/bin_countdown.v#L104-L109}{Fixpoint bin\_cdn\_wkr}` (f : N -> N) (a n : N) (b : nat) : N :=
%  match b with O => 0 | S b' =>
%    if (n <=? a) then 0 else 1 + bin_cdn_wkr f a (f n) b'
%  end.
%\end{lstlisting}
%
%\smallskip
%
%\pause 
%\textbf{Observation.} Recursive argument is budget $b$, \\which should still be in \texttt{nat}.
%
%\bigskip
%
%\textbf{Issue.} $b$ cannot be $n$ due to slow binary $\rightarrow$ unary conversion.
%
%\bigskip
%
%\pause 
%\textbf{Solution.} Use $b \approx \log_2n$ and contractions that contract ``fast enough''. 
%\emph{i.e.} Functions that halve their inputs.
%\end{frame}


\begin{frame}[fragile]
\frametitle{Countdown in Binary}

Our inverse Ackermann implementation runs in $O(n)$ time for unary $n$. \\
What about binary input? \imppar{\textbf{Goal:} $\bigO \big(\log_2n\big)$}.

\bigskip
\pause 
\textbf{Translating Countdown.} Recursive argument for worker is budget $b$, \\which should still be in \texttt{nat}.

\bigskip

\textbf{Issue.} $b$ cannot be $O(n)$ due to slow binary $\rightarrow$ unary conversion.

\bigskip

\pause 
\textbf{Solution.} Use $b \approx \log_2n$ and contractions that contract ``fast enough''. 
\emph{i.e.} Functions that halve their inputs.
\end{frame}



\begin{frame}[fragile]
\frametitle{Binary Contractions and Countdown}

\textbf{Binary contractions.} $f$ is a \textcolor{red}{\emph{binary contraction}} strict above $a$ if $\forall n, f(n)\le n$ and $f(n)\le \lfloor \frac{n+ a}{2} \rfloor$ when $n > a$.

\smallskip

\textbf{Observation.} $f$ shrinks $n$ past $a$ within $\lfloor\log_2(n - a) \rfloor + 1$ steps.

\bigskip

\pause 
New Countdown computation:
% Linked by A
\begin{lstlisting}
`\href{https://github.com/inv-ack/inv-ack/blob/7270e64a2600b771f2b1b1b151f7d13fb2ae6c97/bin_countdown.v#L104-L109}{Fixpoint bin\_cdn\_wkr}` (f : N -> N) (a n : N) (b : nat) : N :=
  match b with O => 0 | S b' =>
    if (n <=? a) then 0 else 1 + bin_cdn_wkr f a (f n) b'
  end.

`\href{https://github.com/inv-ack/inv-ack/blob/7270e64a2600b771f2b1b1b151f7d13fb2ae6c97/bin_countdown.v#L111-L112}{Definition bin\_countdown\_to}` (f : N -> N) (a n : N) : N :=
bin_cdn_wkr f a n (nat_size (n - a)).
\end{lstlisting}

where \texttt{nat\_size x} computes $\lfloor \log_2x \rfloor + 1$ as a \texttt{nat}.
 
\end{frame}


\subsection{Inverse Ackermann in Binary}

\begin{frame}[fragile]
\frametitle{Translating the $\alpha$ Hierarchy}

Must start with a strict binary contraction.

\smallskip

More hard-coding to skip those that are not fast enough.

\bigskip

\pause 
\begin{lstlisting}
`\href{https://github.com/inv-ack/inv-ack/blob/7270e64a2600b771f2b1b1b151f7d13fb2ae6c97/bin_inv_ack.v#L55-L60}{Fixpoint bin\_alpha}` (m : nat) (x : N) : N :=
match m with
  | 0%nat => x - 1          
  | 1%nat => x - 2
  | 2%nat => N.div2 (x - 2) 
  | 3%nat => N.log2 (x + 2) - 2
  | S m'  => bin_countdown_to (bin_alpha m') 1 (bin_alpha m' x)
end.
\end{lstlisting}

\end{frame}


\begin{frame}[fragile]
\frametitle{Translating the Inverse Ackermann Function}

Worker function:

\begin{lstlisting}
`\href{https://github.com/inv-ack/inv-ack/blob/7270e64a2600b771f2b1b1b151f7d13fb2ae6c97/bin_inv_ack.v#L326-L333}{Fixpoint bin\_inv\_ack\_wkr}` (f : N -> N) (n k : N) (b : nat) : N :=
match b with 0%nat  => k | S b' => if n <=? k then k else
  let g := (bin_countdown_to f 1) in
    bin_inv_ack_wkr (compose g f) (g n) (N.succ k) b'
  end.
\end{lstlisting}

\pause 
Same idea: use logarithmic size budget. More hard-coding.

\pause 
\begin{lstlisting}
`\href{https://github.com/inv-ack/inv-ack/blob/7270e64a2600b771f2b1b1b151f7d13fb2ae6c97/bin_inv_ack.v#L335-L342}{Definition bin\_inv\_ack}` (n : N) : N :=
if      (n <=? 1) then 0
else if (n <=? 3) then 1
else if (n <=? 7) then 2
else let f := (fun x => N.log2 (x + 2) - 2)
       in bin_inv_ack_wkr f (f n) 3 (nat_size n).
\end{lstlisting}

\end{frame}


\begin{frame}
\frametitle{Time Complexity of Binary $\alpha$: A Sketch}

Similar to $\alpha$ on unary inputs, $\runtime_{\alpha}(n) \approx \runtime_{\alpha_k}(n)$ for $k\triangleq \alpha(n)$.

\smallskip

\pause
Manual computation: $\runtime_{\alpha_3}(n)\le 2\log_2n + \log_2\log_2n + 3$.

\smallskip

\pause
Suppose $\runtime_{\alpha_i}(n) \le 2\log_2n + \bigO\left(\log_2^{(2)}n\right)$. \pause Then
\begin{equation*}
\begin{aligned}
& \runtime_{\alpha_{i+1}}(n) \ \lessapprox \ 2\log_2n + \bigO\left(\log_2^{(2)}n\right) + 2\log_2^{(2)}n + \bigO\big(\log_2^{(3)}n\big) + \ldots \\
& \ = 2\log_2n + \bigO\left(\log_2^{(2)}n + \log_2^{(3)}n + \log_2^{(4)}n + \ldots\right)
= 2\log_2n + \bigO\left(\log_2^{(2)}n\right)
\end{aligned}
\end{equation*}
\pause
By induction, \imppar{$\runtime_{\alpha}(n) \approx \runtime_{\alpha_k}(n) \lessapprox 2\log_2n + \bigO\left(\log_2^{(2)}n\right) \lessapprox \Theta(\log_2 n)$}.

\end{frame}

\section{Appendix}

\subsection{Runtime of Countdown}
\begin{frame}
\frametitle{Correctness and Runtime of Countdown}

\def\mycolor{blue}
\textbf{Sum by component in each recursive step.}
\begin{equation*}
\begin{array}{l|llllll}
Step & \multicolumn{2}{l}{Initial \ Arguments} & & \big({\color<9>{\mycolor} f^{(0)}(n)}, & n - 0 \big) & \\[3pt]
\pause 0 & {\color<7>{\mycolor} n - 0 > 0}, & {\color<8>{\mycolor} f^{(0)}(n) > a} & \to & \big({\color<9>{\mycolor} f^{(1)}(n)}, & n - 1 \big) & {\color<10>{\mycolor} + 1} \\[3pt]
\pause 1 & {\color<7>{\mycolor} n - 1 > 0}, & {\color<8>{\mycolor} f^{(1)}(n) > a} & \to & \big({\color<9>{\mycolor} f^{(2)}(n)}, & n - 2 \big) & {\color<10>{\mycolor} + 2} \\[3pt]
\pause \vdots & {\color<7>{\mycolor} \vdots} & {\color<8>{\mycolor} \vdots} & \vdots & {\color<9>{\mycolor} \vdots} & \vdots & {\color<10>{\mycolor} \vdots} \\
k-1 & {\color<7>{\mycolor} n - (k-1) > 0}, & {\color<8>{\mycolor} f^{(k-1)}(n) > a} & \to & {\color<9>{\mycolor} \big(f^{(k)}(n)}, & n - k\big) & {\color<10>{\mycolor} + k} \\[3pt]
\pause k & {\color<7>{\mycolor} n - k \mathrel{?} 0}, & {\color{red} {\color<8>{\mycolor} f^{(k)}(n) \le a}} &\to & \multicolumn{3}{l}{0 \hfill \impinline{$k = \cdt{f}{a}(n)$}} \\[3pt] \hline
\pause SUM & \pause {\color<7>{\mycolor} \Theta(k)} & \pause {\color<8>{\mycolor} \Theta((a+1)k)} \topspace{2pt} \pause & & \multicolumn{2}{c}{{\color<9>{\mycolor} \sum_{i=0}^{k-1}\runtime_f\big(f^{i}(n)\big)}} & \pause {\color<10>{\mycolor} \Theta(k)}
\end{array}
\end{equation*}

%TODO: Increase vertical space after \hline

\pause
Substitute $k = \cdt{f}{a}(n)$:
%\begin{equation*}
\imppar{$\displaystyle
\runtime_{\cdt{f}{a}}(n) = \sum_{i=0}^{\cdt{f}{a}(n)-1}\runtime_f\big(f^{i}(n)\big)
+ \Theta\left((a+1)\cdt{f}{a}(n)\right)
$}
%\end{equation*}
\end{frame}


\subsection{Time Complexity for Unary Inputs: $\bigO(n)$}

\begin{frame}
\frametitle{Runtime of Each $\alpha_i$ (Unary Inputs, Asymptotic Bounds)}

$\displaystyle
\alpha_{i+1} = \cdt{\big(\alpha_i\big)}{1}\circ \alpha_i
\quad \implies \quad \runtime_{\alpha_{i+1}}(n) = \runtime_{\alpha_i}(n) + \runtime_{\cdt{\alpha_i}{1}}\big(\alpha_i(n)\big)
$

\smallskip
\pause 
\imppar{
$\displaystyle
\runtime_{\alpha_{i+1}}(n) = \runtime_{\alpha_i}(n) + 
\sum_{i=0}^{\cdt{\alpha_i}{1}(\alpha_i(n)) - 1} \runtime_{\alpha_i}\left(\alpha_i^{(i+1)}(n)\right) + \Theta\left(\cdt{\alpha_i}{1}(\alpha_i(n)) \right)
$
}

\bigskip

\pause 
Substitute $\cdt{\alpha_i}{1}(\alpha_i(n))$ for $\alpha_{i+1}(n)$:
\bigskip

$ \displaystyle 
\runtime_{\alpha_{i+1}}(n)
 = \runtime_{\alpha_i}(n) + 
 \sum_{i=0}^{\alpha_{i+1}(n) - 1} \runtime_{\alpha_i}\left(\alpha_i^{(i+1)}(n)\right) + \Theta\left(\alpha_{i+1}(n) \right)
$

\pause 
\smallskip
\imppar{
$\displaystyle 
\runtime_{\alpha_{i+1}}(n)
	= \sum_{i=0}^{\alpha_{i+1}(n)} \runtime_{\alpha_i}\left(\alpha_i^{(i)}(n)\right) + \Theta\left(\alpha_{i+1}(n) \right)
$
}

\end{frame}


\begin{frame}
\frametitle{Runtime of Each $\alpha_i$ (Unary Inputs, Precise Bounds)}

Countdown runtime:
\begin{equation*}
\textstyle \runtime_{\cdt{f}{a}}(n) =
\sum_{i=0}^{\cdt{f}{a}(n) - 1} \runtime_f\left(f^{(i)}(n)\right)
+ (a + 2)\cdt{f}{a}(n) + f^{\left(\cdt{f}{a}(n)\right)}(n) + 1
\end{equation*}

\smallskip

\pause 
$\alpha_2$ and $\alpha_3$ runtime: $\runtime_{\alpha_2}(n) \le 2n - 2$ and $\runtime_{\alpha_3}(n) \le 4n + 4$.

\bigskip 

\pause 
$\alpha_i$ runtime: $
\textstyle \runtime_{\alpha_{i+1}}(n) \le \sum_{k=0}^{\alpha_{i+1}(n)}\runtime_{\alpha_i}\left( \alpha_i^{(k)}(n)\right) + 3\alpha_{i+1}(n) + 3$.

\bigskip

\pause 
\imppar{
\begin{minipage}{\linewidth}
	\textbf{Theorem.} $\forall i.~\runtime_{\alpha_i}(n) \le 4n + \left(19\cdot 2^{i-3} - 2i - 13\right)\log_2n + 2i = \bigO(n)$, when using $\alpha_1 \triangleq \lambda n. (n-2)$.
\end{minipage}}

\end{frame}



\begin{frame}
\frametitle{Runtime of Inverse Ackermann for Unary Inputs}

\def\mycolor{blue}
\begin{equation*}
\begin{array}{l|llllllll}
Step & \multicolumn{2}{l}{Initial \ Arguments} & & \big(\alpha_1, & {\color<4>{\mycolor} \alpha_1(n)}, & 1, & b - 1 \big) \\[3pt]
1 & {\color<2>{\mycolor} b - 1 > 0}, & {\color<3>{\mycolor} \alpha_1(n) > 1} & \to & \big(\alpha_2, & {\color<4>{\mycolor} \alpha_2(n)}, & 2, & b - 2\big) \\[3pt]
2 & {\color<2>{\mycolor} b - 2 > 0}, & {\color<3>{\mycolor} \alpha_2(n) > 2} & \to & \big(\alpha_3, & {\color<4>{\mycolor} \alpha_3(n)}, & 3, & b - 3\big) \\
\vdots & {\color<2>{\mycolor} \vdots} & {\color<3>{\mycolor} \vdots} & \vdots & \vdots & {\color<5>{\mycolor} \vdots} & \vdots \\
%b > k-1, \ \alpha_{k-1}(n) > k - 1 & \to & \big(\alpha_k, & \alpha_k(n), & k, & b - k\big) & k-1 \\
k & {\color<2>{\mycolor} b - k > 0}, & {\color{red} {\color<3>{\mycolor} \alpha_{k}(n) \le k}} & \to & \multicolumn{4}{l}{k = \alpha(n)} \\[3pt] \hline
SUM & \pause {\color<2>{\mycolor} \Theta(k)} & \pause {\color<3>{\mycolor} \Theta\big(\sum_{i=1}^{k}i\big)} \pause & & \multicolumn{4}{l}{{\color<4>{\mycolor} \sum_{i=1}^{k-1} \runtime_{\cdt{\alpha_i}{1}} (\alpha_i(n))}} \topspace{2pt} \\[5pt]
\pause \color<5>{\mycolor} = & \Theta(k) & \color<5>{\mycolor} \Theta(k^2) & & \multicolumn{4}{l}{\color<5>{\mycolor} \sum_{i=1}^{k-1} \big(\runtime_{\alpha_{i+1}}(n) - \runtime_{\alpha_i}(n) \big)}
\end{array}
\end{equation*}

\smallskip

\pause
Therefore, 
\imppar{$\displaystyle
\begin{aligned}
\runtime_{\alpha}(n)
& = \runtime_{\alpha_{\alpha(n)}}(n) - \runtime_{\alpha_1}(n) + \Theta\left(\alpha(n)^2\right) + \runtime_{\alpha_1}(n) \\
\pause & = \bigO \left(n + 2^{\alpha(n)}\log_2n + \alpha(n)^2 \right) = \bigO (n)
\end{aligned}
$}

\end{frame}

\subsection{Time Complexity for Binary Inputs: $\bigO (\log_2n)$}


\begin{frame}
\frametitle{Runtime of Each $\alpha_i$ (Binary Inputs, Precise Bounds)}

Countdown runtime:
\begin{equation*}
\runtime_{\cdt{f}{a}}(n) \le \sum_{i=0}^{\cdt{f}{a}(n) - 1} \hspace{-6pt}
\runtime_f\big(f^{(i)}(n)\big) \ + \ (\log_2a + 3)\left(\cdt{f}{a}(n) + 1\right) \ + \ 2\log_2n \ + \ \log_2\cdt{f}{a}(n)
\end{equation*}

\pause 
$\alpha_3$ runtime: $\runtime_{\alpha_3}(n) \le 2\log_2n + \log_2\log_2n + 3$.

\bigskip

\pause 
$\alpha_i$ runtime:
$
%\begin{equation*}
\runtime_{\alpha_{i+1}}(n) \le
\sum_{k=0}^{\alpha_{i+1}(n)}\runtime_{\alpha_i}\big(\alpha_i^{(k)}(n)\big)
+ 6\log_2\log_2n + 3.
%\end{equation*}
$

\bigskip

\pause 
\textbf{Theorem.} $\forall i. \ \runtime_{\alpha_i}(n) \le 2\log_2n + \left(3\cdot 2^i - 3i - 13\right)\log_2\log_2n + 3i = \bigO(\log_2n)$, when hard-coding up to $\alpha_3$.

\end{frame}



\begin{frame}
\frametitle{Runtime of Inverse Ackermann for Binary Inputs}

\def\mycolor{blue}
\begin{equation*}
\begin{array}{l|llllllll}
Step & \multicolumn{2}{l}{Initial \ Arguments} & & \big(\alpha_3, & \color<4>{\mycolor} \alpha_3(n), & b - 0, & \color<5>{\mycolor} 3 \big) \\[3pt]
1 & \color<2>{\mycolor} b - 0 > 0, & \color<3>{\mycolor} \alpha_3(n) > 3 & \to & \big(\alpha_4, & \color<4>{\mycolor} \alpha_4(n), & b - 1, & \color<5>{\mycolor} 4 \big) \\[3pt]
2 & \color<2>{\mycolor} b - 1 > 0, & \color<3>{\mycolor} \alpha_4(n) > 4 & \to & \big(\alpha_5, & \color<4>{\mycolor} \alpha_5(n), & b - 2, & \color<5>{\mycolor} 5 \big) \\
\vdots & \color<2>{\mycolor} \vdots & \color<3>{\mycolor} \vdots & \vdots & \vdots & \color<4>{\mycolor} \vdots & \vdots & \color<5>{\mycolor} \vdots \\
%b > k-1, \ \alpha_{k-1}(n) > k - 1 & \to & \big(\alpha_k, & \alpha_k(n), & k, & b - k\big) & k-1 \\
k & \color<2>{\mycolor} b - k > 0, & {\color{red} {\color<3>{\mycolor} \alpha_{k}(n)} \le k} & \to & \multicolumn{4}{l}{k = \alpha(n)} \\[3pt] \hline
SUM & \pause \color<2>{\mycolor} \Theta(k) & \pause \color<3>{\mycolor} \Theta\big(\sum_{i=1}^{k}\log_2i\big) \pause & & \multicolumn{3}{r}{\color<4>{\mycolor} \sum_{i=1}^{k-1} \runtime_{\cdt{\alpha_i}{1}} (\alpha_i(n))} \pause & \color<5>{\mycolor} \Theta(k) \topspace{2pt} \\[5pt]
\pause \color<5>{\mycolor} = & \Theta(k) & \color<5>{\mycolor} \Theta(k\log_2k) & \multicolumn{4}{r}{\color<5>{\mycolor} \sum_{i=1}^{k-1} \big(\runtime_{\alpha_{i+1}}(n) - \runtime_{\alpha_i}(n) \big)} & \Theta(k)
\end{array}
\end{equation*}

\smallskip

\pause 
\imppar{
$\displaystyle
\begin{aligned}
\runtime_{\alpha}(n)
& = \runtime_{\alpha_{\alpha(n)}}(n) - \runtime_{\alpha_3}(n) + \Theta\left(\alpha(n)\log_2 \alpha(n)\right) + \runtime_{\alpha_3}(n) \\
\pause & = \bigO \left(\log_2n + 2^{\alpha(n)}\log_2\log_2n + \alpha(n)\log_2\alpha(n) \right) = \bigO (\log_2n)
\end{aligned}
$}

\end{frame}



%\begin{frame}[allowframebreaks]
%\frametitle{References}
%\bibliographystyle{amsalpha}
%\bibliography{inv_ack}
%\end{frame}

\end{document}